\documentclass[11pt]{article}

\usepackage[margin=1in]{geometry}
\usepackage{amsmath,amssymb,mathtools,bm}
\usepackage{booktabs,array,graphicx,microtype}
\usepackage[colorlinks=true,linkcolor=blue,urlcolor=blue,citecolor=blue]{hyperref}
\usepackage{enumitem}

\newcommand{\ii}{\mathrm{i}}
\newcommand{\ee}{\mathrm{e}}
\newcommand{\RePart}{\operatorname{Re}}
\newcommand{\ImPart}{\operatorname{Im}}
\newcommand{\T}{\mathsf T}
\newcommand{\Mu}{\bm{\mu}}
\newcommand{\cQ}{\mathcal Q}
\newcommand{\cC}{\mathcal C}
\newcommand{\cL}{\mathcal L}
\newcommand{\ip}[2]{\left\langle #1,#2\right\rangle}
\newcolumntype{L}[1]{>{\raggedright\arraybackslash}p{#1}}

\title{Amplitude Selection, V1 Connectivity, and Numerical Tests\\
for the Spatio-temporal Hallucination Model}
\author{}
\date{August 17, 2026}

\begin{document}
\maketitle

\section{Main conclusions}

The two critical waves have amplitudes $A$ and $B$.  Without forcing their
cubic equations are
\begin{equation}
 \dot A=A(a+b|A|^2+c|B|^2),\qquad
 \dot B=B(a+b|B|^2+c|A|^2).
 \label{eq:unforced}
\end{equation}
Here and below the dot denotes the slow time.  There are only three generic
amplitude states: rest, a traveling wave, and a standing wave.  Their
existence and stability follow from two $2\times2$ Jacobians and are derived
in Section~\ref{sec:equilibria}.

Suppose each separated forcing is either absent or has order
$\epsilon^2$, with temporal frequency $2\omega_0$ and spatial wavenumber
$2k_0$.  A direct projection and a parity argument show that no additive
critical forcing is produced: the reduced coefficient previously denoted by
$f$ is identically zero in this regime.  After choosing the temporal and
spatial origins, the two remaining parametric strengths are real and
nonnegative:
\begin{align}
 \dot A&=A(a+b|A|^2+c|B|^2)+\eta_t\overline B+\eta_xB,\notag\\
 \dot B&=B(a+b|B|^2+c|A|^2)+\eta_t\overline A+\eta_xA.
 \label{eq:forced}
\end{align}
Section~\ref{sec:all-forced} derives \eqref{eq:forced} from the forcing
orders, reconstructs every locked phase from the two amplitudes, reduces the
equal-amplitude branches to a quadratic in the squared amplitude, reduces
every unequal-amplitude branch to one scalar equation, and gives exact
stability tests for the rest state, every equal branch, and every remaining
complex equilibrium.

For $K$ cell types, let $\Mu=(\mu_{ij})$ be the signed Gaussian masses and
$\Sigma=(\sigma_{ij})$ the widths.  The exact Fourier connectivity is
$\widehat J(q)_{ij}=\mu_{ij}\exp(-\sigma_{ij}^2|q|^2/2)$.  Each critical wave
has one complex amplitude and one $K$-component eigenvector describing how
all cell types participate.  For any active set $S$ of waves, the generic
cubic equilibrium is $x_S=-G_{SS}^{-1}a_S$, with stability given by one active
radial matrix and the invasion rates of the inactive waves.  Section
\ref{sec:K-framework} derives these formulas and the resonant generalization.
Section~\ref{sec:mixed-domain} extends them to
$\mathbb S^1\times\mathbb R^d$, with one discrete orientation harmonic and a
continuous $d$-dimensional cortical wavevector.

V1 is not biologically an E/I two-population system.  The Allen VISp/ALM
taxonomy has 133 transcriptomic types, and even the deliberately coarse
data-driven V1 model of Billeh et al. uses 17 layer-by-class populations.
MICrONS resolves cell-to-cell wiring at far finer scale, but no complete
physiological gain matrix exists between all transcriptomic types.  Therefore
this report separates measured type identity and connection probability from
calibrated Wilson--Cowan gains.  The full published 17-by-17 probability table
is supplied with the code; unknown entries are not relabeled as measurements.

The computed coefficients for that example are
\begin{equation}
 \alpha=0.08100+0.05273\ii,\qquad
 b=-0.11231-0.15236\ii,\qquad
 c=-1.56376-0.65446\ii.
 \label{eq:numerical-coefficients}
\end{equation}
Since $\RePart c<\RePart b<0$, the unforced cubic theory selects one traveling
wave.  At $a=0.01$ and $\eta_t=\eta_x=0.04$, the real-coefficient projection
predicts a stable standing state.  The full complex calculation places every
equal-amplitude three-population branch outside stability, but the complete
unequal-amplitude quartic reveals a stable exchange-broken locked pair.  This
distinction is invisible if either the phases or the unequal-amplitude
branches are discarded.

The 2025 Drinnenberg et al. preprint identifies a functional split inside
deep-L2/3 SST cells: general ensemble-response (SST-GER) neurons integrate
broad excitatory activity and differ from local SST and PV cells.  Splitting
the field into
$(E,PV,\mathrm{SST\!-\!GER},\mathrm{SST\!-\!local},VIP)$ retains broad
excitation and gives a finite-wave Hopf point at mode 4 and $14.04$ Hz, with
\begin{equation}
 b_5=-1.10320+0.06808\ii,\qquad
 c_5=-2.86511-0.52395\ii.
 \label{eq:five-coefficients}
\end{equation}
Thus the broader functional partition also selects a traveling wave before
forcing.  Under the same symmetric forcing it has a stable complex locked
standing state of amplitude $0.14949$ and phase $172.622^\circ$ modulo
$180^\circ$.  These are existence calculations, not proof that the same
mechanism generates human hallucinations.

On a translation-invariant infinite line, continuous optimization over the
wavevector gives a sharper distinction.  The three-population traveling wave
saturates locally but fails the longitudinal Benjamin--Feir sideband test.
The five-population wave passes finite-wave onset, local invasion, and the
one-dimensional sideband test.  In an isotropic plane or hyperspace,
rotational degeneracy still requires an angular or higher-order transverse
stability calculation before the pattern can be called fully stable.

\section{All amplitude equilibria and their stability}
\label{sec:equilibria}

\subsection{Reduction to the two real amplitudes}

Write
\begin{equation}
 A=r_+\ee^{\ii\theta_+},\qquad
 B=r_-\ee^{\ii\theta_-},\qquad r_+,r_-\geq0,
\end{equation}
and put $a_r=\RePart a$, $b_r=\RePart b$, and $c_r=\RePart c$.  Dividing the
first equation in \eqref{eq:unforced} by $\ee^{\ii\theta_+}$ and separating
real and imaginary parts gives
\begin{align}
 \dot r_+&=r_+(a_r+b_rr_+^2+c_rr_-^2),
 &\dot\theta_+&=\ImPart a+\ImPart b\,r_+^2+\ImPart c\,r_-^2.
 \label{eq:polarA}
\end{align}
The second equation gives the reflected pair
\begin{align}
 \dot r_-&=r_-(a_r+b_rr_-^2+c_rr_+^2),
 &\dot\theta_-&=\ImPart a+\ImPart b\,r_-^2+\ImPart c\,r_+^2.
 \label{eq:polarB}
\end{align}
Thus the imaginary parts change phase and frequency but do not change the
amplitude equilibria.

For compactness, set $x=r_+^2$ and $y=r_-^2$.  Multiplication of
\eqref{eq:polarA} and \eqref{eq:polarB} by $2r_+$ and $2r_-$ gives
\begin{equation}
 \dot x=2x(a_r+b_rx+c_ry),\qquad
 \dot y=2y(a_r+b_ry+c_rx).
 \label{eq:squared}
\end{equation}
At an equilibrium, each right-hand side is zero.  If both $x$ and $y$ are
nonzero, subtraction of the two bracketed equations gives
\begin{equation}
 (b_r-c_r)(x-y)=0.
\end{equation}
Therefore, when $b_r\ne c_r$, $b_r\ne0$, and $b_r+c_r\ne0$, every
equilibrium is one of the following.
The rest state is
\begin{equation}
 (x,y)=(0,0).
 \label{eq:rest}
\end{equation}
The two traveling-wave states are
\begin{equation}
 (x,y)\in
 \left\{\left(-\frac{a_r}{b_r},0\right),
           \left(0,-\frac{a_r}{b_r}\right)\right\}.
 \label{eq:tw}
\end{equation}
The standing-wave state is
\begin{equation}
 (x,y)=\left(-\frac{a_r}{b_r+c_r},
             -\frac{a_r}{b_r+c_r}\right).
 \label{eq:sw}
\end{equation}
The nonzero state exists exactly when the displayed squared amplitude is
positive.  The divisions above lose the following degenerate cases.  A
one-sided state with $x>0$ exists when
\begin{equation}
 a_r+b_rx=0.
 \label{eq:unforced-one-sided-undivided}
\end{equation}
Thus $b_r=0$ gives no one-sided branch when $a_r\ne0$, but gives every
$x>0$ when $a_r=0$.  When $b_r\ne c_r$, a two-sided state must have $x=y>0$
and obey
\begin{equation}
 a_r+(b_r+c_r)x=0.
 \label{eq:unforced-standing-undivided}
\end{equation}
Thus $b_r+c_r=0$ gives no standing branch when $a_r\ne0$, but leaves every
$x=y>0$ unresolved at cubic order when $a_r=0$.  Finally, if $b_r=c_r\ne0$,
the positive two-sided states form the continuum
\begin{equation}
 x+y=-\frac{a_r}{b_r}>0.
 \label{eq:unforced-degenerate-circle}
\end{equation}
If $b_r=c_r=0$, positive two-sided states exist only when $a_r=0$, in which
case every $x,y>0$ is a radial equilibrium.  Each continuum has a zero
radial eigenvalue and requires higher-order terms for selection.

\subsection{Jacobian calculation}

The radial Jacobian before evaluating an equilibrium is
\begin{equation}
 J(r_+,r_-)=
 \begin{pmatrix}
 a_r+3b_rr_+^2+c_rr_-^2 & 2c_rr_+r_-\\
 2c_rr_+r_- & a_r+3b_rr_-^2+c_rr_+^2
 \end{pmatrix}.
 \label{eq:radial-J}
\end{equation}
At rest, both radial eigenvalues are $a_r$.  Hence rest is stable for
$a_r<0$ and loses stability at $a_r=0$.

For a traveling wave $(r_+,r_-)=(\rho,0)$,
$\rho^2=-a_r/b_r$.  Substitution into \eqref{eq:radial-J} gives
\begin{equation}
 J_{\rm TW}=
 \begin{pmatrix}
 -2a_r&0\\0&a_r+c_r\rho^2
 \end{pmatrix},
\end{equation}
and hence
\begin{equation}
 \lambda_{\rm radial}=-2a_r,\qquad
 \lambda_{\rm opposite}=a_r\frac{b_r-c_r}{b_r}.
 \label{eq:tw-eigs}
\end{equation}
The branch can be stable at cubic order only on the supercritical side
$a_r>0$, $b_r<0$.  There the second eigenvalue is negative precisely when
\begin{equation}
 c_r<b_r.
 \label{eq:tw-stability}
\end{equation}

For a standing wave $r_+=r_-=\rho$,
$\rho^2=-a_r/(b_r+c_r)$, and
\begin{equation}
 J_{\rm SW}=2\rho^2
 \begin{pmatrix}b_r&c_r\\c_r&b_r\end{pmatrix}.
\end{equation}
The symmetric and antisymmetric eigenvectors $(1,1)^T$ and $(1,-1)^T$ give
\begin{equation}
 \lambda_{\rm symmetric}=-2a_r,\qquad
 \lambda_{\rm imbalance}=2\rho^2(b_r-c_r).
 \label{eq:sw-eigs}
\end{equation}
Thus a supercritical standing wave requires $a_r>0$, $b_r+c_r<0$, and is
stable to amplitude imbalance precisely when
\begin{equation}
 c_r>b_r.
 \label{eq:sw-stability}
\end{equation}
The traveling- and standing-wave inequalities are complementary.  The two
phase eigenvalues are zero because unforced time translation and spatial
translation act on $\theta_++\theta_-$ and $\theta_+-\theta_-$.  Stability
here therefore means orbital stability modulo those symmetries.

The following table summarizes the nondegenerate supercritical cases.

\begin{center}
\begin{tabular}{@{}llll@{}}
\toprule
State & Squared amplitude & Supercritical existence & Amplitude stability\\
\midrule
Rest & $0$ & all coefficients & $a_r<0$\\
Traveling & $-a_r/b_r$ & $a_r>0$, $b_r<0$ & $c_r<b_r$\\
Standing & $-a_r/(b_r+c_r)$ & $a_r>0$, $b_r+c_r<0$ & $c_r>b_r$\\
\bottomrule
\end{tabular}
\end{center}

\section{Pure \texorpdfstring{$2{:}1$}{2:1} parametric forcing:
\texorpdfstring{$f=0$}{f=0}}
\label{sec:all-forced}

\subsection{Why the reduced additive term is exactly absent}

Use the weakly nonlinear scaling
\begin{equation}
 u=\epsilon u_1+\epsilon^2u_2+\epsilon^3u_3+O(\epsilon^4),
 \qquad T=\epsilon^2t.
 \label{eq:pure21-scaling}
\end{equation}
Each separated input is now restricted to one of two choices:
\begin{center}
\begin{tabular}{@{}lll@{}}
\toprule
Input & Present & Absent\\
\midrule
Temporal &
$\epsilon^2(P\ee^{2\ii\omega_0t}+\overline P\ee^{-2\ii\omega_0t})$ &
$0$\\
Spatial &
$\epsilon^2(S\ee^{2\ii k_0x}+\overline S\ee^{-2\ii k_0x})$ &
$0$\\
\bottomrule
\end{tabular}
\end{center}
Thus $n=2$ when the temporal input is present and $m=2$ when the spatial
input is present; \emph{absent} may equivalently be denoted by an infinite
order.  The label $2{:}1$ is a resonance ratio, not an amplitude ratio: the
forcing frequency or wavenumber is twice the corresponding critical carrier,
\begin{equation}
 \omega_F=2\omega_0,\qquad k_F=2k_0.
 \label{eq:pure21-definition}
\end{equation}
The critical first-order field is
\begin{equation}
 u_1=A(T)\Phi\ee^{\ii(\omega_0t+k_0x)}
 +B(T)\Phi\ee^{\ii(\omega_0t-k_0x)}+\mathrm{c.c.}
 \label{eq:pure21-critical-field}
\end{equation}
The two forcing responses in $u_2$ are
\begin{align}
 p_2&=(2\ii\omega_0I-M_0(0))^{-1}P,\notag\\
 s_2^+&=-M_0(2k_0)^{-1}S,\qquad
 s_2^-=-M_0(-2k_0)^{-1}\overline S .
 \label{eq:pure21-responses}
\end{align}
At order $\epsilon^3$, the quadratic term
$2\cQ(u_1,u_2)$ contains
\begin{align}
 (2\omega_0,0)+(-\omega_0,k_0)&=(\omega_0,k_0),\notag\\
 (2\omega_0,0)+(-\omega_0,-k_0)&=(\omega_0,-k_0),
 \label{eq:pure21-temporal-selection}\\
 (0,2k_0)+(\omega_0,-k_0)&=(\omega_0,k_0),\notag\\
 (0,-2k_0)+(\omega_0,k_0)&=(\omega_0,-k_0).
 \label{eq:pure21-spatial-selection}
\end{align}
Projection by the adjoint critical eigenvector therefore gives
\begin{align}
 \eta_t&=2\Psi^*\cQ(\overline\Phi,p_2),\notag\\
 \eta_x^+&=2\Psi^*\cQ(\Phi,s_2^+),\qquad
 \eta_x^-=2\Psi^*\cQ(\Phi,s_2^-).
 \label{eq:pure21-eta}
\end{align}
Reflection symmetry and a choice of temporal and spatial origins reduce
these to $\eta_t,\eta_x\geq0$ in \eqref{eq:forced}.

The symbol $f$ was not a coefficient of the original unforced amplitude
system.  It was introduced only as shorthand for a possible
amplitude-independent critical residual in a more general forced reduction.
If such a third-order residual had critical coefficients $R_+$ and $R_-$,
the solvability projection would give
\begin{equation}
 f_+=\Psi^*R_+,
 \qquad
 f_-=\Psi^*R_-.
 \label{eq:pure21-f-origin}
\end{equation}
The calculation below shows that $R_+=R_-=0$ for the present exact
$2{:}1$-or-absent inputs, so neither $f_+$ nor $f_-$ appears in
\eqref{eq:forced}.

There are two separate reasons that no additive term occurs. First, a
temporal forcing harmonic has pair $(2\omega_0,0)$ and a spatial forcing
harmonic has pair $(0,\pm2k_0)$; neither equals either critical pair
$(\omega_0,\pm k_0)$. Second, their first mixed product has order
\begin{equation}
 \epsilon^n\epsilon^m=\epsilon^4
 \label{eq:pure21-order-four}
\end{equation}
and pair $(2\omega_0,\pm2k_0)$, so it is neither third order nor critical.
More strongly, every amplitude-independent monomial formed only from the
doubled forcing harmonics has pair
\begin{equation}
 (2j\omega_0,2\ell k_0),\qquad j,\ell\in\mathbb Z,
 \label{eq:pure21-even-lattice}
\end{equation}
which cannot equal $(\omega_0,\pm k_0)$. Hence an analytic nonlinearity
cannot generate an additive critical term from exact $2{:}1$ separated
forcing at any order. In this report's regime,
\begin{equation}
 \boxed{f=0.}
 \label{eq:pure21-fzero}
\end{equation}
An additive term would require a $1{:}1$ spatio-temporal harmonic, leakage
outside the doubled forcing lattice, or explicit translation-breaking
heterogeneity; none is assumed here.

If the temporal frequency is
$2\omega_0+\epsilon^2\delta_t$, the parametric terms initially contain
$\ee^{\ii\delta_tT}$. The transformation
\begin{equation}
 A=\widetilde A\ee^{\ii\delta_tT/2},\qquad
 B=\widetilde B\ee^{\ii\delta_tT/2}
 \label{eq:pure21-detuning-transform}
\end{equation}
makes the system autonomous and replaces $a$ by
$a-\ii\delta_t/2$. We therefore allow $a,b,c\in\mathbb C$ below and use
\eqref{eq:forced} without tildes.

\subsection{Exact polar equations and the definition of locking}

Write
\begin{equation}
 A=r\ee^{\ii\phi},\qquad B=\rho\ee^{\ii\psi},
 \qquad r,\rho\geq0,
 \label{eq:general-polar-form}
\end{equation}
and separate
\begin{align}
 a+br^2+c\rho^2&=\alpha_A+\ii\beta_A,&
 a+b\rho^2+cr^2&=\alpha_B+\ii\beta_B.
 \label{eq:polar-local-coefficients}
\end{align}
For $r,\rho>0$, multiplication of \eqref{eq:forced} by
$\ee^{-\ii\phi}$ and $\ee^{-\ii\psi}$ gives exactly
\begin{align}
 \dot r={}&r\alpha_A+\eta_t\rho\cos(\phi+\psi)
 +\eta_x\rho\cos(\phi-\psi),
 \label{eq:polar-r}\\
 r\dot\phi={}&r\beta_A-\eta_t\rho\sin(\phi+\psi)
 -\eta_x\rho\sin(\phi-\psi),
 \label{eq:polar-phi}\\
 \dot\rho={}&\rho\alpha_B+\eta_t r\cos(\phi+\psi)
 +\eta_xr\cos(\phi-\psi),
 \label{eq:polar-rho}\\
 \rho\dot\psi={}&\rho\beta_B-\eta_t r\sin(\phi+\psi)
 +\eta_xr\sin(\phi-\psi).
 \label{eq:polar-psi}
\end{align}
A phase-locked equilibrium satisfies
\begin{equation}
 \dot r=\dot\rho=\dot\phi=\dot\psi=0.
 \label{eq:general-locking-definition}
\end{equation}
The physical carrier still oscillates at frequency $\omega_0$; locking means
that both slow amplitudes and both phase offsets relative to the forcing are
constant.

Put
\begin{equation}
 \Theta=\phi+\psi,\qquad \Delta=\phi-\psi.
 \label{eq:phase-sum-difference}
\end{equation}
The radial equations contain
\begin{equation}
 C=\eta_t\cos\Theta+\eta_x\cos\Delta
 \label{eq:phase-C}
\end{equation}
and are equivalent to
\begin{equation}
 r\alpha_A+\rho C=0,\qquad
 \rho\alpha_B+rC=0.
 \label{eq:general-radial-locking}
\end{equation}
Consequently
\begin{equation}
 C=-\frac{r}{\rho}\alpha_A
   =-\frac{\rho}{r}\alpha_B,\qquad
 r^2\alpha_A=\rho^2\alpha_B.
 \label{eq:general-radial-compatibility}
\end{equation}
The phase equations give
\begin{align}
 T:=\eta_t\sin\Theta
 &=\frac12\left(\frac{r}{\rho}\beta_A
                 +\frac{\rho}{r}\beta_B\right),
 \label{eq:phase-T}\\
 X:=\eta_x\sin\Delta
 &=\frac12\left(\frac{r}{\rho}\beta_A
                 -\frac{\rho}{r}\beta_B\right).
 \label{eq:phase-X}
\end{align}
For each $\varepsilon_t,\varepsilon_x\in\{-1,1\}$, positive amplitudes give
a locked state if and only if
\begin{align}
 |T|&\leq\eta_t,\qquad |X|\leq\eta_x,
 \label{eq:phase-sine-bounds}\\
 -\frac{r}{\rho}\alpha_A
 &=\varepsilon_t\sqrt{\eta_t^2-T^2}
  +\varepsilon_x\sqrt{\eta_x^2-X^2},
 \label{eq:general-phase-amplitude-closure}
\end{align}
together with the last equality in
\eqref{eq:general-radial-compatibility}. The phases are then
\begin{align}
 \Theta&=\operatorname{atan2}\!\left(
 T,\varepsilon_t\sqrt{\eta_t^2-T^2}\right),\notag\\
 \Delta&=\operatorname{atan2}\!\left(
 X,\varepsilon_x\sqrt{\eta_x^2-X^2}\right),
 \label{eq:general-phase-reconstruction}\\
 \phi&=\frac{\Theta+\Delta}{2},\qquad
 \psi=\frac{\Theta-\Delta}{2}.
 \label{eq:individual-phase-reconstruction}
\end{align}
These conditions are necessary and sufficient, not merely necessary.
If $\eta_t=0$, the first sine condition is replaced by $T=0$, the term
$\eta_t\cos\Theta$ is zero, and $\Theta$ is a neutral phase; the first
$\operatorname{atan2}$ formula is then omitted.  If $\eta_x=0$, the second
sine condition is replaced by $X=0$, the term $\eta_x\cos\Delta$ is zero,
and $\Delta$ is neutral; the second $\operatorname{atan2}$ formula is
omitted.  After assigning arbitrary values to the neutral phases,
\eqref{eq:individual-phase-reconstruction} still reconstructs $\phi$ and
$\psi$.  Thus no undefined $\operatorname{atan2}(0,0)$ is used when an input
is absent.

\subsection{A factorization that exhausts the positive amplitudes}

Let $a_r=\RePart a$, $b_r=\RePart b$, and $c_r=\RePart c$. Expanding the
last condition in \eqref{eq:general-radial-compatibility} gives
\begin{align}
 0={}&r^2(a_r+b_rr^2+c_r\rho^2)
 -\rho^2(a_r+b_r\rho^2+c_rr^2)\notag\\
 ={}&(r^2-\rho^2)\left[a_r+b_r(r^2+\rho^2)\right].
 \label{eq:forced-amplitude-factorization}
\end{align}
Therefore every positive equilibrium satisfies
\begin{equation}
 \boxed{r=\rho}
\end{equation}
or, when $b_r\ne0$, it satisfies
\begin{equation}
 \boxed{r^2+\rho^2=-\frac{a_r}{b_r}}.
 \label{eq:equal-or-unequal}
\end{equation}
If $b_r=0$ and $a_r\ne0$, no unequal-amplitude state exists. If
$a_r=b_r=0$, the factorization is degenerate and the necessary-and-sufficient
closure \eqref{eq:phase-sine-bounds}--
\eqref{eq:general-phase-amplitude-closure} is used directly.

For the unequal class define only
\begin{equation}
 R=r^2+\rho^2=-\frac{a_r}{b_r},\qquad
 \Pi=r\rho,\qquad
 D=r^2-\rho^2=\pm\sqrt{R^2-4\Pi^2},
 \label{eq:unequal-invariants}
\end{equation}
where $R>0$ and $0<\Pi<R/2$. The quantities in the phase closure reduce to
\begin{align}
 C&=(b_r-c_r)\Pi,
 \label{eq:unequal-C}\\
 T(\Pi)&=\frac{a_iR+b_iR^2+2(c_i-b_i)\Pi^2}{2\Pi},
 \label{eq:unequal-T}\\
 X(\Pi,D)&=\frac{D(a_i+b_iR)}{2\Pi},
 \label{eq:unequal-X}
\end{align}
where $a_i=\ImPart a$, $b_i=\ImPart b$, and $c_i=\ImPart c$. Hence all
unequal-amplitude equilibria are the roots $0<\Pi<R/2$ of
\begin{equation}
 (b_r-c_r)\Pi
 =\varepsilon_t\sqrt{\eta_t^2-T(\Pi)^2}
 +\varepsilon_x\sqrt{\eta_x^2-X(\Pi,D)^2},
 \label{eq:unequal-scalar-equation}
\end{equation}
subject to $|T(\Pi)|\leq\eta_t$ and
$|X(\Pi,D)|\leq\eta_x$. The two signs of $D$ give the exchanged states.
Then
\begin{equation}
 r^2=\frac{R+D}{2},\qquad \rho^2=\frac{R-D}{2},
 \label{eq:unequal-reconstruction}
\end{equation}
and the phases follow from
\eqref{eq:general-phase-reconstruction}--
\eqref{eq:individual-phase-reconstruction}.

The remaining scalar problem can be reduced to a quartic without losing the
sign information.  Put $z=\Pi^2$, $d_r=b_r-c_r$, and
\begin{equation}
 L=a_i+b_iR,
 \qquad
 N(z)=a_iR+b_iR^2+2(c_i-b_i)z.
 \label{eq:unequal-polynomial-shortcuts}
\end{equation}
The two squared radical arguments, after multiplication by $4z$, are
\begin{align}
 U_0(z)&=4\eta_t^2z-N(z)^2,
 \label{eq:unequal-U0}\\
 V_0(z)&=4\eta_x^2z-(R^2-4z)L^2.
 \label{eq:unequal-V0}
\end{align}
Thus the unsquared equation has
\begin{equation}
 C=d_r\sqrt z,
 \qquad
 \eta_t^2-T^2=\frac{U_0(z)}{4z},
 \qquad
 \eta_x^2-X^2=\frac{V_0(z)}{4z}.
 \label{eq:unequal-radical-arguments}
\end{equation}
Eliminating the two square roots gives the explicit quartic equation
\begin{equation}
 16d_r^2z^2V_0(z)
 -\left[4d_r^2z^2+V_0(z)-U_0(z)\right]^2=0.
 \label{eq:unequal-quartic}
\end{equation}
Every unequal equilibrium gives a root of \eqref{eq:unequal-quartic}.  A
quartic root is retained exactly when
\begin{equation}
 0<z<\frac{R^2}{4},
 \qquad U_0(z)\geq0,
 \qquad V_0(z)\geq0,
 \label{eq:unequal-quartic-filter}
\end{equation}
and it satisfies the original signed equation
\eqref{eq:unequal-scalar-equation} for at least one pair
$(\varepsilon_t,\varepsilon_x)$.  This final unsquared test removes every
extraneous root introduced by squaring.  Consequently there are at most four
admissible unequal products $\Pi$ at generic cubic order, and the two signs
of $D$ reconstruct the exchanged states.

\subsection{Closed equal-amplitude solutions at every phase}

Let $r=\rho>0$ and first assume $\eta_x>0$. Then
\eqref{eq:phase-X} gives $\sin\Delta=0$, so
\begin{equation}
 B=\varepsilon A,\qquad \varepsilon\in\{1,-1\},\qquad
 A=r\ee^{\ii\phi},\qquad x=r^2.
 \label{eq:equal-general-phase-ansatz}
\end{equation}
With $s=b+c$ and
\begin{equation}
 h_\varepsilon(x)=a+\varepsilon\eta_x+sx,
 \label{eq:equal-h}
\end{equation}
the equilibrium equation is
\begin{equation}
 h_\varepsilon(x)+\varepsilon\eta_t\ee^{-2\ii\phi}=0.
 \label{eq:equal-general-phase-balance}
\end{equation}
For $\eta_t>0$, taking the modulus gives
\begin{equation}
 |s|^2x^2+2\RePart\!\left[
 (a+\varepsilon\eta_x)\overline s\right]x
 +|a+\varepsilon\eta_x|^2-\eta_t^2=0.
 \label{eq:equal-general-phase-quadratic}
\end{equation}
Put
\begin{equation}
 D_\varepsilon=
 \RePart\!\left[(a+\varepsilon\eta_x)\overline s\right]^2
 -|s|^2\left(|a+\varepsilon\eta_x|^2-\eta_t^2\right).
 \label{eq:equal-general-phase-discriminant}
\end{equation}
For $s\ne0$, all squared amplitudes are
\begin{equation}
 x_{\varepsilon,\pm}=
 \frac{-\RePart[(a+\varepsilon\eta_x)\overline s]
 \pm\sqrt{D_\varepsilon}}{|s|^2},
 \label{eq:equal-general-phase-roots}
\end{equation}
retaining exactly the roots with $D_\varepsilon\geq0$ and
$x_{\varepsilon,\pm}>0$. The corresponding phases are
\begin{align}
 \ee^{-2\ii\phi}
 &=-\frac{h_\varepsilon(x)}{\varepsilon\eta_t},
 \label{eq:equal-general-phase-exponential}\\
 \phi&=-\frac12\arg\!\left[
 -\frac{h_\varepsilon(x)}{\varepsilon\eta_t}\right]\pmod{\pi}.
 \label{eq:equal-general-phase-angle}
\end{align}
The two values modulo $2\pi$ differ by $\pi$ and give the sign-related
states $(A,B)$ and $(-A,-B)$.

If $s=0$, the modulus equation is independent of $x$.  There is no positive
equal-amplitude branch unless
\begin{equation}
 |a+\varepsilon\eta_x|=\eta_t.
 \label{eq:equal-degenerate-condition}
\end{equation}
When this equality holds and $\eta_t>0$, every $x>0$ satisfies the
cubic-order modulus condition and its phase is given by
\eqref{eq:equal-general-phase-exponential}.  If $\eta_t=0$, the undivided
condition is $a+\varepsilon\eta_x=0$ and the common phase is neutral.  Each
continuum is a degeneracy of the cubic truncation and must be resolved by
quintic terms.

For temporal forcing alone, $\eta_x=0$ and $\eta_t>0$, the positive
equal-amplitude family
does not require $B=\pm A$.  Instead, for $x=r^2$ it is
\begin{equation}
 A=r\ee^{\ii(\Theta+\Delta)/2},\qquad
 B=r\ee^{\ii(\Theta-\Delta)/2},\qquad
 |a+sx|=\eta_t,\qquad
 \ee^{-\ii\Theta}=-\frac{a+sx}{\eta_t},
 \label{eq:temporal-only-equal-family}
\end{equation}
where $\Delta$ is arbitrary.  This gives the general locked phase rather
than only the special choices $B=\pm A$.

The cases in which one or both inputs are absent follow without a limiting
ambiguity:
\begin{center}
\begin{tabular}{@{}L{0.18\textwidth}L{0.36\textwidth}L{0.34\textwidth}@{}}
\toprule
Forcing present & Equal-amplitude condition & Phase information\\
\midrule
Temporal and spatial &
$|a+\varepsilon\eta_x+sx|=\eta_t$ &
$\Delta=0$ or $\pi$; $\phi$ from
\eqref{eq:equal-general-phase-angle}\\
Temporal only &
$|a+sx|=\eta_t$ &
$\Theta$ is locked by
$a+sx+\eta_t\ee^{-\ii\Theta}=0$; $\Delta$ is neutral\\
Spatial only &
$a+\varepsilon\eta_x+sx=0$ &
$\Delta=0$ or $\pi$; the common phase is neutral\\
Neither &
$a+sx=0$ for a stationary complex equilibrium &
both group phases are neutral; Section~\ref{sec:equilibria} gives the more
general relative equilibria\\
\bottomrule
\end{tabular}
\end{center}
With no forcing, a radial equilibrium whose imaginary frequency correction
is nonzero is a relative equilibrium rather than a stationary point in the
chosen rotating frame.

\subsection{Exceptional one-sided states}

Polar coordinates omit $r=0$ or $\rho=0$, so return to
\eqref{eq:forced}. If $A=0$ and $B=\rho\ee^{\ii\psi}\ne0$, then
\begin{equation}
 \eta_t\overline B+\eta_xB=0,\qquad
 a+b\rho^2=0.
 \label{eq:one-sided-conditions}
\end{equation}
If at least one parametric forcing is present, the first equation implies
\begin{equation}
 \eta_t=\eta_x>0,\qquad
 \ee^{2\ii\psi}=-1,
 \label{eq:one-sided-phase}
\end{equation}
and, when $b\ne0$, the second requires the complex number $-a/b$ to be real
and positive, with $\rho^2=-a/b$.  If $b=0$, it instead requires $a=0$ and
leaves every $\rho>0$ unresolved at cubic order.  Exchanging $A$ and $B$
gives the other one-sided state.
If $\eta_t\ne\eta_x$, no nonzero one-sided equilibrium exists. These
codimension conditions explain why generic parametric forcing converts
unforced traveling waves into two-component locked states.
When $\eta_t=\eta_x=0$, the first condition in
\eqref{eq:one-sided-conditions} disappears and the usual unforced
one-sided traveling-wave relative equilibria of
Section~\ref{sec:equilibria} return.

\subsection{Why the equilibrium list is exhaustive}

Every point $(A,B)\in\mathbb C^2$ belongs to exactly one of the disjoint
amplitude classes
\begin{equation}
 (|A|,|B|)=(0,0),\qquad
 |A||B|=0<|A|+|B|,\qquad
 |A|=|B|>0,\qquad
 |A|\ne|B|,\quad |A||B|>0.
 \label{eq:forced-exhaustive-partition}
\end{equation}
The first class is rest.  The second is exactly the one-sided calculation
\eqref{eq:one-sided-conditions}.  In the third class, the phase equations
give either \eqref{eq:equal-general-phase-roots}--
\eqref{eq:equal-general-phase-angle}, the temporal-only family
\eqref{eq:temporal-only-equal-family}, or one of the absent-input rows in
the table above.  In the fourth class,
\eqref{eq:forced-amplitude-factorization} forces
\eqref{eq:equal-or-unequal}, and the necessary-and-sufficient phase closure
reduces it to \eqref{eq:unequal-scalar-equation}.  Hence these formulas omit
no complex equilibrium.  The exceptional coefficient cases $b_r=0$ or
$a_r=b_r=0$ are not divided by; they are decided directly by
\eqref{eq:general-radial-compatibility} and
\eqref{eq:general-phase-amplitude-closure}.

\subsection{Closed real-coefficient classification}

For real $a,b,c$ and real $A,B$, put
\begin{equation}
 q=\eta_t+\eta_x,\qquad s=b+c,\qquad d=b-c.
 \label{eq:forcing-combinations}
\end{equation}
Adding and subtracting the two equilibrium equations gives
\begin{align}
 0&=(A-B)\left[a-q+b(A^2+B^2)+dAB\right],
 \label{eq:forced-difference}\\
 0&=(A+B)\left[a+q+b(A^2+B^2)-dAB\right].
 \label{eq:forced-sum}
\end{align}
Therefore the complete real list is:
\begin{align}
 &(A,B)=(0,0),\notag\\
 &(A,B)=(r,r),\qquad r^2=-\frac{a+q}{s}>0,
 \label{eq:real-in-phase}\\
 &(A,B)=(r,-r),\qquad r^2=-\frac{a-q}{s}>0,
 \label{eq:real-antiphase}
\end{align}
plus exchange-broken states. When $bd\ne0$, the latter satisfy
\begin{equation}
 A^2+B^2=-\frac{a}{b},\qquad AB=\frac{q}{d}.
 \label{eq:real-broken-invariants}
\end{equation}
Writing $U=-a/b$ and $P=q/d$, they exist exactly when
\begin{equation}
 U+2P\geq0,\qquad U-2P\geq0,
 \label{eq:real-broken-existence}
\end{equation}
and all sign-related and exchanged roots are
\begin{equation}
 A=\frac{\sigma\sqrt{U+2P}+\tau\sqrt{U-2P}}{2},\qquad
 B=\frac{\sigma\sqrt{U+2P}-\tau\sqrt{U-2P}}{2},
 \quad \sigma,\tau\in\{-1,1\}.
 \label{eq:real-broken-roots}
\end{equation}
Equalities in \eqref{eq:real-broken-existence} merge this family with an
equal or anti-equal branch. If $b=0$, $d=0$, or $s=0$, the undivided
factorization \eqref{eq:forced-difference}--\eqref{eq:forced-sum} gives the
corresponding degenerate cases without losing a branch.

\subsection{Exact stability: rest and every equal-amplitude branch}

At the rest state, exchange-symmetric and exchange-antisymmetric
perturbations satisfy
\begin{equation}
 \dot z_\kappa=(a+\kappa\eta_x)z_\kappa
 +\kappa\eta_t\overline z_\kappa,\qquad \kappa\in\{1,-1\}.
 \label{eq:rest-exchange-modes}
\end{equation}
Their characteristic polynomials and eigenvalues are
\begin{align}
 0={}&\lambda^2-2\RePart(a+\kappa\eta_x)\lambda
 +|a+\kappa\eta_x|^2-\eta_t^2,
 \label{eq:rest-characteristic}\\
 \lambda_{\kappa,\pm}
 ={}&\RePart a+\kappa\eta_x
 \pm\sqrt{\eta_t^2-(\ImPart a)^2}.
 \label{eq:rest-eigenvalues}
\end{align}
The square root is imaginary when its argument is negative. The rest state
is asymptotically stable if and only if, for both $\kappa=\pm1$,
\begin{equation}
 \RePart(a+\kappa\eta_x)<0,\qquad
 |a+\kappa\eta_x|^2>\eta_t^2.
 \label{eq:rest-stability}
\end{equation}

Now let $B=\varepsilon A$, $A=r\ee^{\ii\phi}$, and $x=r^2$. The full
$4\times4$ linearization splits exactly into the two exchange sectors
$\delta B=\kappa\delta A$. In sector $\kappa$ it is the scalar
conjugate-linear equation
\begin{equation}
 \delta\dot A=P_\kappa\delta A+Q_\kappa\delta\overline A,
 \label{eq:equal-sector-linearization}
\end{equation}
where
\begin{align}
 P_\kappa&=
 a+\left[2b+c+\kappa\varepsilon c\right]x+\kappa\eta_x,
 \label{eq:equal-sector-P}\\
 Q_\kappa&=
 x\ee^{2\ii\phi}(b+\kappa\varepsilon c)+\kappa\eta_t.
 \label{eq:equal-sector-Q}
\end{align}
The two eigenvalues in that sector obey
\begin{equation}
 \lambda^2-2\RePart P_\kappa\,\lambda
 +|P_\kappa|^2-|Q_\kappa|^2=0
 \label{eq:equal-sector-characteristic}
\end{equation}
and therefore
\begin{equation}
 \lambda_{\kappa,\pm}
 =\RePart P_\kappa
 \pm\sqrt{|Q_\kappa|^2-(\ImPart P_\kappa)^2}.
 \label{eq:equal-sector-eigenvalues}
\end{equation}
Thus an equal-amplitude locked state is asymptotically stable if and only if,
for each $\kappa=\pm1$,
\begin{equation}
 \boxed{\RePart P_\kappa<0,\qquad
 |P_\kappa|^2>|Q_\kappa|^2.}
 \label{eq:equal-sector-stability}
\end{equation}
If one forcing is absent, the continuous symmetry it leaves unbroken
produces the expected zero eigenvalue; stability is then orbital after
removing that neutral group direction.

\subsection{Exact stability at every remaining complex root}

Let the right-hand side of the first equation in \eqref{eq:forced} be
$F_A$. Its Wirtinger derivatives are
\begin{align}
 \partial_A F_A&=a+2b|A|^2+c|B|^2,
 &\partial_{\overline A}F_A&=bA^2,\notag\\
 \partial_B F_A&=cA\overline B+\eta_x,
 &\partial_{\overline B}F_A&=cAB+\eta_t.
 \label{eq:wirtinger}
\end{align}
The derivatives of $F_B$ follow by exchanging $A$ and $B$. The real matrix
of $z\mapsto Pz+Q\overline z$ is
\begin{equation}
 \mathcal R(P,Q)=
 \begin{pmatrix}
 \RePart(P+Q)&-\ImPart(P-Q)\\
 \ImPart(P+Q)& \RePart(P-Q)
 \end{pmatrix}.
 \label{eq:realification}
\end{equation}
Using \eqref{eq:realification} on the four derivative pairs in
\eqref{eq:wirtinger} and its exchanged copy produces the exact real
$4\times4$ Jacobian
\begin{equation}
 J=\begin{pmatrix}
 \mathcal R(a+2b|A|^2+c|B|^2,bA^2)&
 \mathcal R(cA\overline B+\eta_x,cAB+\eta_t)\\
 \mathcal R(cB\overline A+\eta_x,cAB+\eta_t)&
 \mathcal R(a+2b|B|^2+c|A|^2,bB^2)
 \end{pmatrix}.
 \label{eq:full-forced-Jacobian}
\end{equation}
This formula applies at any root reconstructed above, including the
codimension-one-sided states.

For a coefficient-only stability test, write
\begin{equation}
 \det(\lambda I-J)=
 \lambda^4+h_1\lambda^3+h_2\lambda^2+h_3\lambda+h_4,
 \label{eq:quartic-characteristic}
\end{equation}
where no eigenvalue calculation is required:
\begin{align}
 h_1&=-\operatorname{tr}J,\notag\\
 h_2&=\frac12\left[(\operatorname{tr}J)^2-\operatorname{tr}(J^2)\right],
 \notag\\
 h_3&=-\frac16\left[(\operatorname{tr}J)^3
 -3\operatorname{tr}J\,\operatorname{tr}(J^2)
 +2\operatorname{tr}(J^3)\right],\notag\\
 h_4&=\det J.
 \label{eq:quartic-coefficients}
\end{align}
The root is asymptotically stable if and only if the quartic
Routh--Hurwitz inequalities hold:
\begin{equation}
 h_1,h_2,h_3,h_4>0,\qquad
 h_1h_2>h_3,\qquad
 h_1h_2h_3>h_3^2+h_1^2h_4.
 \label{eq:quartic-hurwitz}
\end{equation}
Equations \eqref{eq:unequal-scalar-equation} and
\eqref{eq:quartic-hurwitz} therefore give a complete existence-and-stability
algorithm for every unequal complex phase-locked state.  Applying the same
Jacobian and Hurwitz test to \eqref{eq:one-sided-conditions} gives the exact
stability of every exceptional one-sided state.

\subsection{Application to the computed three- and five-population coefficients}

For the real-coefficient projection, take $a=0.01$ and
$\eta_t=\eta_x=0.04$. The real parts of
\eqref{eq:numerical-coefficients} give
\begin{equation}
 (A,B)=(0,0),\qquad (A,B)=(r,r),\qquad(A,B)=(-r,-r),
 \qquad r=0.23173.
\end{equation}
The rest state is unstable. Each nonzero state has eigenvalues
\begin{equation}
 (-0.18000,-0.00412,-0.08000,-0.08000)
\end{equation}
and is asymptotically stable within this real-coefficient projection. These
real states are not equilibria of the full complex-coefficient normal form.

At zero detuning, retain the complete complex values of $b,c$ and apply
\eqref{eq:equal-general-phase-roots}--
\eqref{eq:equal-general-phase-angle}. The exact sector formula
\eqref{eq:equal-sector-eigenvalues} gives:
\begin{center}
\begin{tabular}{@{}lrrrrl@{}}
\toprule
Model & $\varepsilon$ & $r$ & $\phi$ (degrees) & $\max\RePart\lambda$ & Stability\\
\midrule
Three-population & $+1$ & 0.20564 & 150.732 & 0.00100 & unstable\\
Three-population & $+1$ & 0.07843 &  93.563 & 0.07839 & unstable\\
Three-population & $-1$ & 0.07617 &  86.639 & 0.08000 & unstable\\
Five-population  & $+1$ & 0.14949 & 172.622 & $-0.07861$ & stable\\
Five-population  & $+1$ & 0.05024 &  90.824 & 0.07991 & unstable\\
Five-population  & $-1$ & 0.05016 &  89.178 & 0.08000 & unstable\\
\bottomrule
\end{tabular}
\end{center}
Thus all equal-amplitude branches of the full three-population model are
unstable at this forcing, although the large branch is only weakly unstable.
The five-population model has the stable complex locked state
\begin{equation}
 (A,B)=\left(0.14949\ee^{\ii\,172.622^\circ},
              0.14949\ee^{\ii\,172.622^\circ}\right)
 \label{eq:five-complex-locked-state}
\end{equation}
and its sign-related partner.

The equal-amplitude table is not the complete numerical list.  For both
models the one-sided condition fails because $-a/b$ is not real.  Solving
the quartic \eqref{eq:unequal-quartic} and applying the unsquared filter
\eqref{eq:unequal-scalar-equation} gives the following representatives:
\begin{center}
\begin{tabular}{@{}lrrrrrl@{}}
\toprule
Model & $r$ & $\rho$ & $\phi$ (degrees) & $\psi$ (degrees) &
$\max\RePart\lambda$ & Stability\\
\midrule
Three-population & 0.06978 & 0.29011 & $-24.559$ & $-66.184$ &
 0.01390 & unstable\\
Three-population & 0.15984 & 0.25197 & $-26.250$ & $-35.442$ &
 $-0.00181$ & stable\\
Five-population & 0.00077 & 0.09520 & $-16.561$ & 89.904 &
 0.07233 & unstable\\
\bottomrule
\end{tabular}
\end{center}
Each row has an exchanged partner obtained by interchanging
$(r,\phi)$ with $(\rho,\psi)$, and each of those has a sign-related partner
obtained by adding $180^\circ$ to both phases.  Substitution into the full
four-real-dimensional system gives residuals below $2\times10^{-17}$ before
rounding the displayed values.

Therefore the complete calculation changes the activation-pattern
conclusion for the three-population coefficients: although every
equal-amplitude locked branch is unstable, there is a stable
exchange-broken phase-locked pair.  For the five-population coefficients,
the stable state remains the equal-amplitude branch
\eqref{eq:five-complex-locked-state}; its unequal pair is unstable.

\begin{figure}[ht]
\centering
\includegraphics[width=\textwidth]{output/simulations/amplitude_phase_diagrams.png}
\caption{Amplitude phase diagrams using the computed real cubic coefficients.
White circles are stable states and the cross is the unstable unforced
standing wave.}
\label{fig:phase}
\end{figure}


\section{Closed K-population Gaussian and wave framework}
\label{sec:K-framework}

\subsection{The two connectivity matrices}

Let $u=(u_1,\ldots,u_K)^T$ collect the activities of $K$ cell populations.
The signed Gaussian mass and width matrices are
\begin{equation}
 \Mu=(\mu_{ij})\in\mathbb R^{K\times K},\qquad
 \Sigma=(\sigma_{ij})\in(0,\infty)^{K\times K},
\end{equation}
where row $i$ is the target and column $j$ is the source.  Normalize the
Gaussian to unit mass and set
\begin{align}
 G_\sigma(x)&=\frac{1}{(2\pi\sigma^2)^{d/2}}
 \exp\!\left(-\frac{|x|^2}{2\sigma^2}\right),
 &J_{ij}(x)&=\mu_{ij}G_{\sigma_{ij}}(x).
 \label{eq:gaussian}
\end{align}
Thus $\mu_{ij}=\int J_{ij}$ is the total signed connection gain, not its peak.
If data give a Gaussian peak $p_{ij}$, then
$\mu_{ij}=p_{ij}(2\pi)^{d/2}\sigma_{ij}^{d}$.  The field equation is
\begin{equation}
 \T\partial_tu=-u+F(J*u+I),\qquad
 \T=\operatorname{diag}(\tau_1,\ldots,\tau_K),
 \label{eq:field}
\end{equation}
with $F$ acting componentwise.  On $\mathbb R^d$, or on a torus using wrapped
Gaussians, its Fourier connectivity matrix is the closed expression
\begin{equation}
 \widehat J(q)=\Mu\odot
 \exp\!\left[-\frac{|q|^2}{2}\Sigma^{\odot2}\right],
 \qquad
 \widehat J(q)_{ij}=\mu_{ij}\ee^{-\sigma_{ij}^2|q|^2/2},
 \label{eq:gaussian-symbol}
\end{equation}
where the exponential and powers are entrywise.  This is the only place where
the Gaussian assumption enters the spectral calculation.

\subsection{Homogeneous equilibria and exact field stability}

Every homogeneous equilibrium $u^*$ is a root of the $K$ equations
\begin{equation}
 u^*=F(\Mu u^*+I).
 \label{eq:homogeneous}
\end{equation}
This is an exact finite-dimensional reduction because every kernel has unit
mass.  No formula by radicals exists for an arbitrary choice of $K$ nonlinear
functions $F_i$.  On an affine or threshold-linear branch
$F(h)=Sh+c$, however, the equilibrium is explicitly
\begin{equation}
 u^*=(I-S\Mu)^{-1}(SI+c),
 \label{eq:affine-equilibrium}
\end{equation}
provided $I-S\Mu$ is nonsingular and the result satisfies the inequalities
defining that branch.  For an invertible sigmoid, \eqref{eq:homogeneous} is
equivalently $F^{-1}(u^*)-\Mu u^*=I$ and is solved as a $K$-variable root.

Let $D_r=\operatorname{diag}(F_i^{(r)}(h_i^*))$, where
$h^*=\Mu u^*+I$.  A perturbation $v\ee^{\lambda t+\ii q\cdot x}$ satisfies
\begin{equation}
 \lambda v=M(q)v,\qquad
 M(q)=\T^{-1}\left[-I+D_1\widehat J(q)\right].
 \label{eq:Mq}
\end{equation}
Consequently the equilibrium is linearly asymptotically stable precisely when
\begin{equation}
 \sup_{q\in\mathcal Q}\max_{1\leq s\leq K}
 \RePart\lambda_s(M(q))<0,
 \label{eq:exact-stability}
\end{equation}
where $\mathcal Q$ is continuous on $\mathbb R^d$ and is the reciprocal
lattice on a periodic domain.  A zero of the left side at $q=0$ is a
homogeneous instability; a zero first reached at $|q|>0$ is a pattern-forming
instability.  This criterion includes every cell type and every spatial wave,
without reducing $\Mu$ to an effective E/I difference.

For modal vectors $v,w$ at wavevectors $q_1,q_2$, write
$\ell_q(v)=\widehat J(q)v$.  Taylor expansion of $F$ gives the symmetric
quadratic and cubic maps
\begin{align}
 \cQ_{q_1,q_2}(v,w)
 &=\frac12\T^{-1}D_2
 \left[\ell_{q_1}(v)\odot\ell_{q_2}(w)\right],
 \label{eq:Qgaussian}\\
 \cC_{q_1,q_2,q_3}(u,v,w)
 &=\frac16\T^{-1}D_3
 \left[\ell_{q_1}(u)\odot\ell_{q_2}(v)\odot\ell_{q_3}(w)\right].
 \label{eq:Cgaussian}
\end{align}
Equations \eqref{eq:gaussian-symbol}, \eqref{eq:Mq},
\eqref{eq:Qgaussian}, and \eqref{eq:Cgaussian} express the complete cubic
local dynamics in terms of $\Mu$, $\Sigma$, $\T$, and the derivatives of the
$K$ firing-rate functions.

\subsection{Cell types carried by critical waves}

Suppose $m$ critical wave types have wavevectors $q_\alpha$, frequencies
$\omega_\alpha$, and right and adjoint cell-type vectors satisfying
\begin{equation}
 M(q_\alpha)\Phi_\alpha=\ii\omega_\alpha\Phi_\alpha,
 \qquad M(q_\alpha)^*\Psi_\alpha=-\ii\omega_\alpha\Psi_\alpha,
 \qquad \Psi_\alpha^*\Phi_\alpha=1.
 \label{eq:eigenvectors}
\end{equation}
The leading wave superposition is
\begin{equation}
 u_i(x,t)=u_i^*+2\varepsilon\RePart\sum_{\alpha=1}^m
 A_\alpha(T)\Phi_{i\alpha}
 \ee^{\ii(q_\alpha\cdot x+\omega_\alpha t)}+O(\varepsilon^2),
 \qquad T=\varepsilon^2t.
 \label{eq:cell-wave-superposition}
\end{equation}
There is one complex dynamical variable $A_\alpha$ per wave, not per cell
type.  The entry $\Phi_{i\alpha}$ gives the amplitude and phase with which
cell type $i$ participates in wave $\alpha$.  Thus the $K$-cell biology is
retained inside each wave rather than replaced by $K m$ redundant amplitude
variables.

For completeness, the general cubic projection can be written without
listing every conjugate case.  Give every critical wave a signed copy
$r=(\alpha,s)$, $s=\pm1$, and define
\begin{equation}
 z_{\alpha,+}=A_\alpha,\quad z_{\alpha,-}=\overline A_\alpha,
 \quad q_{\alpha,s}=s q_\alpha,\quad
 \omega_{\alpha,s}=s\omega_\alpha,\quad
 \phi_{\alpha,+}=\Phi_\alpha,\quad
 \phi_{\alpha,-}=\overline\Phi_\alpha.
\end{equation}
For a noncritical pair $(r,s)$, its slaved second-order cell-type response is
\begin{equation}
 H_{rs}=\left[\ii(\omega_r+\omega_s)I-M(q_r+q_s)\right]^{-1}
 \cQ_{q_r,q_s}(\phi_r,\phi_s).
 \label{eq:general-slaved-mode}
\end{equation}
The inverse is taken only off the critical eigenspace.  Summing over ordered
signed indices, the amplitude equation through cubic order is
\begin{align}
 \dot A_\alpha={}&\lambda_\alpha A_\alpha
 +\!\sum_{\substack{q_r+q_s=q_\alpha\\
                     \omega_r+\omega_s=\omega_\alpha}}
 \Psi_\alpha^*\cQ_{q_r,q_s}(\phi_r,\phi_s)z_rz_s\notag\\
 &+\!\sum_{\substack{q_r+q_s+q_t=q_\alpha\\
                     \omega_r+\omega_s+\omega_t=\omega_\alpha}}
 \Psi_\alpha^*\!\left[
 \cC_{q_r,q_s,q_t}(\phi_r,\phi_s,\phi_t)
 +2\cQ_{q_r,q_s+q_t}(\phi_r,H_{st})\right]z_rz_sz_t
 +\sum_{\beta=1}^m\left(
 \eta^x_{\alpha\beta}A_\beta
 +\eta^t_{\alpha\beta}\overline A_\beta\right).
 \label{eq:general-resonant-amplitudes}
\end{align}
Here $\lambda_\alpha$ is the projected parameter detuning.
$\eta^x_{\alpha\beta}$ is nonzero only when an order-$2$ spatial forcing
harmonic scatters wave $\beta$ into wave $\alpha$, while
$\eta^t_{\alpha\beta}$ is nonzero
only when an order-$2$ temporal forcing harmonic scatters the conjugate of
wave $\beta$ into wave $\alpha$.  They are the multi-wave analogues of
$\eta_x$ and $\eta_t$ in \eqref{eq:forced}.  There is no additive vector in
the order-$2$ or absent regime by \eqref{eq:pure21-even-lattice}.  The two
equalities under each nonlinear sum are conservation of wavevector and
temporal frequency.  If a quadratic pair is itself critical, it is retained
in the first sum and is not inverted in
\eqref{eq:general-slaved-mode}.  Equation
\eqref{eq:general-resonant-amplitudes} is the desired wave-based dynamical
system for arbitrary $K$ and any finite collection of interacting waves.

\subsection{Closed equilibria and stability of a wave superposition}

Away from special quadratic and phase resonances, translation and time-shift
symmetry reduce \eqref{eq:general-resonant-amplitudes} to
\begin{equation}
 \dot A_\alpha=A_\alpha\left(\lambda_\alpha+
 \sum_{\beta=1}^m g_{\alpha\beta}|A_\beta|^2\right),
 \qquad \alpha=1,\ldots,m.
 \label{eq:m-wave-amplitudes}
\end{equation}
Put $x_\alpha=|A_\alpha|^2$, $a_\alpha=\RePart\lambda_\alpha$, and
$G_{\alpha\beta}=\RePart g_{\alpha\beta}$.  Then
\begin{equation}
 \dot x_\alpha=2x_\alpha\left(a_\alpha+
 \sum_{\beta=1}^mG_{\alpha\beta}x_\beta\right).
 \label{eq:intensity-system}
\end{equation}
Every equilibrium is indexed by its active wave set $S$.  If $G_{SS}$ is
nonsingular, its closed form is
\begin{equation}
 x_S=-G_{SS}^{-1}a_S>0,\qquad x_\alpha=0\quad(\alpha\notin S).
 \label{eq:all-wave-equilibria}
\end{equation}
Thus all generic cubic wave equilibria are obtained by checking the at most
$2^m$ supports; negative components are discarded.  At such a state the
active radial Jacobian and the inactive invasion rates are
\begin{equation}
 J_S=2\operatorname{diag}(x_S)G_{SS},\qquad
 \nu_\alpha=2\left(a_\alpha+G_{\alpha S}x_S\right)
 \quad(\alpha\notin S).
 \label{eq:wave-stability}
\end{equation}
For one active wave $\alpha$, this reads
$x_\alpha=-a_\alpha/G_{\alpha\alpha}$.  For two active waves $\alpha,\beta$,
put
$\Delta=G_{\alpha\alpha}G_{\beta\beta}
-G_{\alpha\beta}G_{\beta\alpha}$.  Their explicit intensities are
\begin{equation}
 x_\alpha=\frac{G_{\alpha\beta}a_\beta
 -G_{\beta\beta}a_\alpha}{\Delta},\qquad
 x_\beta=\frac{G_{\beta\alpha}a_\alpha
 -G_{\alpha\alpha}a_\beta}{\Delta}.
 \label{eq:two-wave-general-root}
\end{equation}
These are retained only when both are positive; their radial and invasion
stability then follows directly from \eqref{eq:wave-stability}.
The wave superposition is asymptotically stable modulo phase symmetries if
every eigenvalue of $J_S$ has negative real part and every $\nu_\alpha<0$.
Its slow phase velocities are
\begin{equation}
 \dot\theta_\alpha=\ImPart\lambda_\alpha+
 \sum_\beta\ImPart g_{\alpha\beta}x_\beta.
\end{equation}
They give nonlinear frequency corrections; forcing or resonant monomials can
lock them and remove the neutral phase directions.  When such terms are
present, equilibria are roots of the explicit polynomial system
\eqref{eq:general-resonant-amplitudes}, and stability is decided by its real
$2m\times2m$ Jacobian.  There is no universal radical formula for that fully
resonant problem, but the dynamical system and Jacobian remain closed.

\subsection{The counter-propagating pair as a special case}

At a single critical $q_0$, write $\Phi=\Phi_1$, $\Psi=\Psi_1$ and
\begin{equation}
 R(\Omega,q)=\left(\ii\Omega I-M(q)\right)^{-1}.
\end{equation}
The second-order fields required for the self coefficient are
\begin{align}
 H_{11}^0&=2R(0,0)\cQ_{q_0,-q_0}(\Phi,\overline\Phi),\notag\\
 H_{20}^{2q_0}&=R(2\omega_0,2q_0)\cQ_{q_0,q_0}(\Phi,\Phi),
\end{align}
and the self coefficient is
\begin{align}
 b=\Psi^*\big[&2\cQ_{q_0,0}(\Phi,H_{11}^0)
 +2\cQ_{-q_0,2q_0}(\overline\Phi,H_{20}^{2q_0})\notag\\
 &+3\cC_{q_0,q_0,-q_0}(\Phi,\Phi,\overline\Phi)\big].
 \label{eq:bformula}
\end{align}
For the counter-propagating coefficient, define
\begin{align}
 H_{20}^{0}&=2R(2\omega_0,0)\cQ_{q_0,-q_0}(\Phi,\Phi),\notag\\
 H_{11}^{2q_0}&=2R(0,2q_0)\cQ_{q_0,q_0}(\Phi,\overline\Phi),
\end{align}
and
\begin{align}
 c=\Psi^*\big[&2\cQ_{q_0,0}(\Phi,H_{11}^0)
 +2\cQ_{q_0,0}(\overline\Phi,H_{20}^{0})
 +2\cQ_{-q_0,2q_0}(\Phi,H_{11}^{2q_0})\notag\\
 &+6\cC_{q_0,-q_0,q_0}(\Phi,\Phi,\overline\Phi)\big].
 \label{eq:cformula}
\end{align}
If a scalar bifurcation parameter changes $M$ by $M_1$, then
\begin{equation}
 \alpha=\Psi^*M_1(q_0)\Phi,\qquad a=\mu\alpha.
 \label{eq:alphaformula}
\end{equation}
This recovers \eqref{eq:unforced}: $g_{11}=g_{22}=b$ and
$g_{12}=g_{21}=c$.  Every coefficient is obtained from $K$-vectors,
$K\times K$ matrix solves, and adjoint projections.

\section{Rigorous translation-invariant domain extensions}

\subsection{Extension 1: the infinite line}
\label{sec:nonperiodic}

First let $\Omega=\mathbb R$.  The Fourier transform is continuous rather
than a Fourier series.  Equations \eqref{eq:gaussian-symbol} and
\eqref{eq:Mq} become
\begin{align}
 \widehat J(q)_{ij}&=\mu_{ij}\ee^{-\sigma_{ij}^2q^2/2},
 &q&\in\mathbb R,\label{eq:R-symbol}\\
 \partial_t\widehat v(q,t)&=M(q)\widehat v(q,t),
 &M(q)&=\T^{-1}\left[-I+D_1\widehat J(q)\right].
 \label{eq:R-linearization}
\end{align}
Equivalently, the linearized operator is the direct integral
\begin{equation}
 \mathcal F\mathcal L\mathcal F^{-1}
 =\int_{\mathbb R}^{\oplus}M(q)\,dq,
 \qquad
 \operatorname{spec}(\mathcal L)
 =\overline{\bigcup_{q\in\mathbb R}\operatorname{spec}(M(q))}.
 \label{eq:R-direct-integral}
\end{equation}
Plancherel's theorem gives
\begin{equation}
 \|v(t)\|_{L^2}^2=
 \int_{\mathbb R}\|\ee^{tM(q)}\widehat v(q,0)\|_{\mathbb C^K}^2\,dq.
\end{equation}
Hence a homogeneous equilibrium is exponentially stable when there are
$C,\gamma>0$ such that
\begin{equation}
 \|\ee^{tM(q)}\|\leq C\ee^{-\gamma t}
 \quad(q\in\mathbb R,t\geq0),
 \label{eq:R-uniform-stability}
\end{equation}
For the Gaussian symbol this is equivalent to the spectral-gap condition
\begin{equation}
 \sup_{q\in\mathbb R}\max_s\RePart\lambda_s(M(q))<0.
 \label{eq:R-spectral-gap}
\end{equation}
To see the nontrivial implication, use
$M(q)\to-\T^{-1}$ as $|q|\to\infty$, so the matrix family is continuous on
the compactification $\mathbb R\cup\{\infty\}$.  If the left side of
\eqref{eq:R-spectral-gap} is $-2\gamma$, choose one contour $\Gamma$ enclosing
all spectra in $\{\RePart z\le-\gamma\}$.  Then
\begin{equation}
 \ee^{tM(q)}=\frac{1}{2\pi\ii}\int_\Gamma
 \ee^{tz}(zI-M(q))^{-1}\,dz,\qquad
 \sup_q\|\ee^{tM(q)}\|\le
 \frac{|\Gamma|}{2\pi}\sup_{q,z\in\Gamma}
 \|(zI-M(q))^{-1}\|\ee^{-\gamma t}.
 \label{eq:R-dunford-bound}
\end{equation}
The reverse implication follows by applying
\eqref{eq:R-uniform-stability} to an eigenvector of $M(q)$.
At threshold, a carrier satisfies
\begin{equation}
 M(q_0)\Phi=\ii\omega_0\Phi.
 \label{eq:R-carrier-eigenproblem}
\end{equation}
It is a generalized Fourier mode rather than an $L^2(\mathbb R)$
eigenfunction, because $\Phi\ee^{\ii q_0x}$ has constant modulus.
Consequently the whole-line center spectrum is not a finite-dimensional
eigenspace.  The two carrier amplitudes describe spatially uniform envelopes
only; localized perturbations require the modulation equation derived in
Extension 2.

The extension used from this point onward is translation invariant.  Thus
the spatial domain is $\mathbb R^d$, not a bounded subset, the coupling is a
true convolution, and every nonlinear interaction retains exact wavevector
conservation.  A bounded truncation is a different model because its edges
break the translation group and pin the phase of a pattern.

\subsection{Extension 2: non-periodic space in arbitrary dimension}
\label{sec:high-dimensional}

Let $x\in\mathbb R^d$ and allow a positive-definite covariance
$C_{ij}\in\mathbb R^{d\times d}$ for each ordered cell-type pair:
\begin{align}
 G_{C_{ij}}(x)&=\frac{
 \exp[-\tfrac12x^TC_{ij}^{-1}x]}
 {(2\pi)^{d/2}\det(C_{ij})^{1/2}},
 &J_{ij}(x)&=\mu_{ij}G_{C_{ij}}(x),\label{eq:anisotropic-G}\\
 \widehat J(q)_{ij}&=\mu_{ij}
 \exp[-\tfrac12q^TC_{ij}q],
 &q&\in\mathbb R^d.
 \label{eq:anisotropic-symbol}
\end{align}
The scalar-width model is recovered by $C_{ij}=\sigma_{ij}^2I_d$.  The full
stability criterion is
\begin{equation}
 \sup_{q\in\mathbb R^d}\max_s\RePart\lambda_s\!\left(
 \T^{-1}[-I+D_1\widehat J(q)]\right)<0.
 \label{eq:Rd-stability}
\end{equation}

\subsubsection{Exact linear condition for a hallucination-like pattern}

In this report, a hallucination-like state means a self-sustained,
nonuniform cortical solution of the field model.  This mathematical
definition does not by itself establish a subjective human hallucination.
Let $\delta$ move the homogeneous equilibrium through onset and define
\begin{equation}
 M(q,\delta)=\T^{-1}\left[-I+D_1(\delta)\widehat J(q,\delta)\right],
 \qquad
 s(q,\delta)=\max_{1\le j\le K}
 \RePart\lambda_j(M(q,\delta)).
 \label{eq:pattern-spectral-abscissa}
\end{equation}
The homogeneous state is exponentially stable exactly when
\begin{equation}
 \sup_{q\in\mathbb R^d}s(q,\delta)<0.
 \label{eq:homogeneous-Rd-gap}
\end{equation}
A finite-wave stationary or oscillatory instability occurs at
$\delta=0$ if there are $q_0\ne0$, a simple eigenvalue branch
$\lambda_*$, and $\omega_0\geq0$ such that
\begin{equation}
 \lambda_*(q_0,0)=\ii\omega_0,
 \qquad
 s(0,0)<0,
 \qquad
 \sup_{q\in\mathbb R^d}s(q,0)=0.
 \label{eq:finite-wave-threshold}
\end{equation}
The crossing is nondegenerate at an isolated anisotropic maximum when
\begin{equation}
 \partial_\delta\RePart\lambda_*(q_0,0)>0,
 \qquad
 \nabla_q\RePart\lambda_*(q_0,0)=0,
 \qquad
 \xi^T\nabla_q^2\RePart\lambda_*(q_0,0)\xi<0
 \quad(\xi\ne0).
 \label{eq:finite-wave-nondegeneracy}
\end{equation}
For isotropic kernels, write $r=|q|$.  The critical set is the sphere
$|q|=q_0$, so the tangential Hessian directions are neutral by rotation
symmetry and the correct radial condition is
\begin{equation}
 \partial_r\RePart\lambda_*(q_0,0)=0,
 \qquad
 \partial_r^2\RePart\lambda_*(q_0,0)<0.
 \label{eq:isotropic-radial-maximum}
\end{equation}
The case $\omega_0=0$ is a stationary finite-wave instability.  The case
$\omega_0>0$ is a finite-wave Hopf instability and produces oscillating or
traveling cortical patterns.  Immediately above onset, the necessary linear
pattern condition is therefore
\begin{equation}
 s(0,\delta)<0,
 \qquad
 \sup_{q\ne0}s(q,\delta)>0.
 \label{eq:linear-hallucination-condition}
\end{equation}
This says that a nonzero spatial scale grows while the spatially homogeneous
mode remains damped.  It is not yet sufficient, because nonlinear growth
must saturate and the resulting pattern must resist all long-wave
modulations.

Assume the critical eigenvalue bands are simple and smooth, all other bands
remain uniformly separated from the imaginary axis, and the first crossing
occurs at a finite set $q_1,\ldots,q_m$.  The cell-resolved expansion is then
\begin{equation}
 u_i(x,t)=u_i^*+2\varepsilon\RePart\sum_{\alpha=1}^m
 A_\alpha(X,T)\Phi_{i\alpha}
 \ee^{\ii(q_\alpha\cdot x+\omega_\alpha t)}+O(\varepsilon^2),
 \quad X=\varepsilon x,\ T=\varepsilon^2t.
 \label{eq:Rd-envelope-ansatz}
\end{equation}
Let $\lambda_\alpha(q,\delta)$ be the simple eigenvalue branch and put
$\delta=\varepsilon^2\mu$.  Its Taylor expansion is
\begin{align}
 \lambda_\alpha(q_\alpha+\varepsilon\kappa,\varepsilon^2\mu)
 ={}&\ii\omega_\alpha
 +\varepsilon\nabla_q\lambda_\alpha(q_\alpha,0)\cdot\kappa\notag\\
 &+\varepsilon^2\left[
 \alpha_\alpha\mu+\frac12\kappa^TH_\alpha\kappa\right]
 +O(\varepsilon^3),
 \label{eq:dispersion-expansion}
\end{align}
where
\begin{equation}
 \alpha_\alpha=\partial_\delta\lambda_\alpha(q_\alpha,0),
 \qquad H_\alpha=\nabla_q^2\lambda_\alpha(q_\alpha,0).
\end{equation}
After removing the order-$\varepsilon$ transport by a comoving coordinate,
the envelope equation whose Fourier symbol matches
\eqref{eq:dispersion-expansion} is
\begin{equation}
 \partial_TA_\alpha=\alpha_\alpha\mu A_\alpha
 -\frac12\sum_{r,s=1}^d(H_\alpha)_{rs}
 \partial_{X_rX_s}A_\alpha
 +A_\alpha\sum_{\beta=1}^m
 g_{\alpha\beta}|A_\beta|^2.
 \label{eq:high-dimensional-CGL}
\end{equation}
Indeed, substituting $A_\alpha=\ee^{\ii\kappa\cdot X+\Lambda T}$ into the
linear part gives
\begin{equation}
 \Lambda=\alpha_\alpha\mu+\frac12\kappa^TH_\alpha\kappa,
\end{equation}
which is exactly the order-$\varepsilon^2$ term in
\eqref{eq:dispersion-expansion}.

The finite-dimensional wave equations describe only envelopes that are
constant in $X$.  To test a constant equilibrium against every localized
sideband, denote the local right-hand side of
\eqref{eq:high-dimensional-CGL} by $\mathcal N(A,\overline A)$ after removing
the derivative term.  Define its doubled real-linear Jacobian by
\begin{equation}
 \mathbb J_0=
 D_{(A,\overline A)}
 \begin{pmatrix}\mathcal N\\ \overline{\mathcal N}\end{pmatrix}
 \bigg|_{A=A^*}.
 \label{eq:sideband-J0}
\end{equation}
For a perturbation proportional to
$\ee^{\ii\kappa\cdot X+\Lambda T}$, put
\begin{equation}
 d_\alpha(\kappa)=\frac12\kappa^TH_\alpha\kappa.
 \label{eq:sideband-symbol}
\end{equation}
The exact $2m\times2m$ sideband matrix is
\begin{equation}
 \mathbb J(\kappa)=\mathbb J_0+
 \operatorname{diag}\!\left(
 d_1(\kappa),\ldots,d_m(\kappa),
 \overline{d_1(\kappa)},\ldots,
 \overline{d_m(\kappa)}\right).
 \label{eq:sideband-J}
\end{equation}
Therefore orbital modulational stability requires
\begin{equation}
 \max_{\lambda\in\operatorname{spec}(\mathbb J(\kappa))}
 \RePart\lambda<0
 \quad(\kappa\ne0),
 \qquad
 \operatorname{spec}(\mathbb J(0))
 \subset\{z:\RePart z<0\}\cup\{0\}.
 \label{eq:sideband-stability}
\end{equation}
The kernel at $\kappa=0$ must consist exactly of the phase and translation
directions generated by symmetry.  If
$\lambda^{\rm sym}_\ell(\kappa)$ are the spectral branches issuing from
that kernel, diffusive decay requires constants $c,\kappa_0>0$ such that
\begin{equation}
 \RePart\lambda^{\rm sym}_\ell(\kappa)
 \leq-c|\kappa|^2
 \qquad(0<|\kappa|<\kappa_0).
 \label{eq:sideband-diffusive-decay}
\end{equation}
The Jacobian test in
\eqref{eq:wave-stability} is precisely the $\kappa=0$ part and is not, by
itself, a sideband-stability test.

\subsubsection{Necessary and sufficient conditions within the amplitude regime}

Let $S$ be the active set of a finite planform and let
$x_S=-G_{SS}^{-1}a_S$.  Within the validity of the cubic modulation
equations, existence and homogeneous-amplitude stability are equivalent to
\begin{equation}
 x_S>0,
 \qquad
 \max_{\lambda\in\operatorname{spec}(J_S)}\RePart\lambda<0,
 \qquad
 \nu_\alpha<0\quad(\alpha\notin S).
 \label{eq:local-hallucination-condition}
\end{equation}
Here $J_S$ and $\nu_\alpha$ are given exactly by
\eqref{eq:wave-stability}.  The planform is a stable hallucination-like state
precisely when 
\eqref{eq:local-hallucination-condition},
\eqref{eq:sideband-stability}, and
\eqref{eq:sideband-diffusive-decay} all hold.  A spatially varying prescribed
input is absent from this criterion: it would select an origin and break
translation invariance.  Spatially uniform parameter changes or temporal
forcing can preserve the spatial translation group.

For isotropic kernels, $M(q)$ depends only on $|q|$ and the critical set is
the sphere $|q|=q_0$.  It is then not a finite-dimensional center eigenspace.
The correct leading superposition is
\begin{equation}
 u_i(x,t)=u_i^*+2\varepsilon\RePart
 \int_{\mathbb S^{d-1}}
 A(\xi,X,T)\Phi_i(\xi)
 \ee^{\ii(q_0\xi\cdot x+\omega_0t)}\,d\xi+O(\varepsilon^2).
 \label{eq:spherical-superposition}
\end{equation}
Its cubic interactions obey
\begin{equation}
 q_0(\xi_1+\xi_2-\xi_3)=q_0\xi,\qquad
 \omega_1+\omega_2-\omega_3=\omega,
 \label{eq:spherical-resonance}
\end{equation}
and their coefficients are still obtained from
\eqref{eq:Qgaussian}--\eqref{eq:general-slaved-mode}.  Choosing a finite
planform $\{\xi_1,\ldots,\xi_m\}$ and restricting to spatially constant
envelopes reduces
\eqref{eq:spherical-superposition} to the finite wave system
\eqref{eq:m-wave-amplitudes}.  Its equilibria and homogeneous-amplitude
stability are exactly \eqref{eq:all-wave-equilibria}--
\eqref{eq:wave-stability}; sidebands are governed by
\eqref{eq:sideband-J}.

For two wavevectors separated by $\vartheta$,
\begin{equation}
 |q_j+q_k|=2q_0\cos(\vartheta/2),\qquad
 |q_j-q_k|=2q_0\sin(\vartheta/2),
\end{equation}
so $g_{jk}=c(\vartheta)$ is obtained by evaluating the resolvents at these
two magnitudes.  For a regular $m$-wave planform with common cross-coupling
$c_*$,
\begin{equation}
 \rho^2=-\frac{a_r}{b_r+(m-1)c_{*,r}},\qquad
 \lambda_{\rm sym}=-2a_r,\qquad
 \lambda_{\rm tr}=2\rho^2(b_r-c_{*,r}).
 \label{eq:Rd-regular-planform}
\end{equation}
The transverse eigenvalue $\lambda_{\rm tr}$ has multiplicity $m-1$.
For a single active traveling direction, the local conditions reduce to
\begin{equation}
 \rho^2=-\frac{a_r}{b_r}>0,
 \qquad
 -2a_r<0,
 \qquad
 2a_r\left(1-\frac{c_r}{b_r}\right)<0.
 \label{eq:single-wave-local-condition}
\end{equation}
Thus $a_r>0$, $b_r<0$, and $c_r<b_r$ give a supercritical traveling wave
that is stable against invasion by its counterpropagating partner.  With
$d\geq2$, however, the isotropic critical sphere contains every direction.
If $c_r(\vartheta)$ is the real cross-coupling at angular separation
$\vartheta$, stability of a single active direction against all other
critical directions requires
\begin{equation}
 \nu(\vartheta)=2a_r\left(1-
 \frac{c_r(\vartheta)}{b_r}\right)<0
 \qquad(0<\vartheta\leq\pi).
 \label{eq:angular-invasion-condition}
\end{equation}
Thus the one-dimensional coefficient $c=c(\pi)$ is a complete angular
invasion test on $\mathbb R$, whose critical set is $\{-q_0,q_0\}$, but not
on a plane or hyperspace.  There one must verify
\eqref{eq:angular-invasion-condition} over the full angular continuum.

On the infinite line, the single-wave envelope can be written in a comoving
and corotating frame as
\begin{equation}
 \partial_TA=aA+D\partial_{XX}A+b|A|^2A,
 \qquad
 D=-\frac12\partial_q^2\lambda_*(q_0,0).
 \label{eq:single-wave-CGL}
\end{equation}
Put $D=D_r+\ii D_i$ and $b=b_r+\ii b_i$.  The uniform traveling state is
\begin{equation}
 A_*(T)=\rho\ee^{\ii\Omega T},
 \qquad
 \rho^2=-\frac{a_r}{b_r},
 \qquad
 \Omega=a_i+b_i\rho^2.
 \label{eq:single-wave-CGL-state}
\end{equation}
Write $A=A_*(1+u+\ii v)$ and take a sideband proportional to
$\ee^{\ii\kappa X}$.  Direct linearization gives
\begin{equation}
 \partial_T
 \begin{pmatrix}u\\v\end{pmatrix}
 =\begin{pmatrix}
 2\rho^2b_r-D_r\kappa^2&D_i\kappa^2\\
 2\rho^2b_i-D_i\kappa^2&-D_r\kappa^2
 \end{pmatrix}
 \begin{pmatrix}u\\v\end{pmatrix}.
 \label{eq:single-wave-sideband-matrix}
\end{equation}
Its trace and determinant are
\begin{equation}
 \operatorname{tr}\mathbb J(\kappa)
 =2\rho^2b_r-2D_r\kappa^2,
 \qquad
 \det\mathbb J(\kappa)
 =\kappa^2\left[|D|^2\kappa^2
 -2\rho^2(b_rD_r+b_iD_i)\right].
 \label{eq:single-wave-trace-determinant}
\end{equation}
The two-dimensional Hurwitz criterion is therefore necessary and sufficient
for every nonzero sideband:
\begin{equation}
 D_r>0,
 \qquad
 b_r<0,
 \qquad
 b_rD_r+b_iD_i<0.
 \label{eq:single-wave-BF-condition}
\end{equation}
Equivalently, when $D_rb_r\ne0$, the Benjamin--Feir number must satisfy
\begin{equation}
 \mathcal B=1+\frac{D_i}{D_r}\frac{b_i}{b_r}>0.
 \label{eq:Benjamin-Feir-number}
\end{equation}

The continuous-wavevector calculation changes the periodic-box numbers
slightly.  If $\chi$ is the common coupling scale, the three-population
infinite-line onset is
\begin{equation}
 \begin{aligned}
 \chi_c&=0.9999159,
 &q_0&=0.380160\,(0.1\,{\rm mm})^{-1},\\
 f_0&=14.1042\,{\rm Hz},
 &s(0,0)&=-0.0017022\,{\rm ms}^{-1},\\
 \partial_q^2\RePart\lambda_*&=-0.0642473.
 \end{aligned}
 \label{eq:three-pop-continuous-onset}
\end{equation}
Its continuous-onset coefficients and longitudinal diffusion are
\begin{equation}
 \begin{aligned}
 b_3^\infty&=-0.073616-0.172897\ii,
 &c_3^\infty(\pi)&=-1.485777-0.626915\ii,\\
 D_3&=0.032124-0.024065\ii,
 &\mathcal B_3&=-0.75944.
 \end{aligned}
 \label{eq:three-pop-continuous-coefficients}
\end{equation}
At $a_r=0.01$, its local amplitude is $\rho=0.36856$ and its
counterpropagating invasion rate is $-0.38365$.  It passes the local
amplitude test but fails \eqref{eq:single-wave-BF-condition}; the selected
traveling wave is longitudinally sideband-unstable on the infinite line at
cubic order.

For the five-population model, the corresponding values are
\begin{equation}
 \begin{aligned}
 \chi_c&=1.0390375,
 &q_0&=0.491786\,(0.1\,{\rm mm})^{-1},\\
 f_0&=13.7341\,{\rm Hz},
 &s(0,0)&=-0.0072509\,{\rm ms}^{-1},\\
 \partial_q^2\RePart\lambda_*&=-0.1153424.
 \end{aligned}
 \label{eq:five-pop-continuous-onset}
\end{equation}
and
\begin{equation}
 \begin{aligned}
 b_5^\infty&=-0.778597+0.066322\ii,
 &c_5^\infty(\pi)&=-2.537480-0.379713\ii,\\
 D_5&=0.057671+0.002692\ii,
 &\mathcal B_5&=0.99602.
 \end{aligned}
 \label{eq:five-pop-continuous-coefficients}
\end{equation}
At $a_r=0.01$, its local amplitude is $\rho=0.11333$ and its
counterpropagating invasion rate is $-0.04518$.  It satisfies the complete
one-dimensional cubic test: finite-wave onset, supercritical saturation,
counter-wave exclusion, and longitudinal sideband stability.

A numerical scan of the angular coefficient gives
\begin{equation}
 \max_{0<\vartheta\leq\pi}c_{3,r}^\infty(\vartheta)
 =-0.14723<b_{3,r}^\infty,
 \qquad
 \max_{0<\vartheta\leq\pi}c_{5,r}^\infty(\vartheta)
 =-1.23240<b_{5,r}^\infty.
 \label{eq:angular-scan-results}
\end{equation}
Thus both models pass the computed angular invasion test, although only the
five-population model also passes the longitudinal sideband test.  This scan
is numerical evidence, not an analytic bound over the angular continuum.
For an isotropic plane or hyperspace, rotational degeneracy makes the real
quadratic curvature vanish in directions tangent to the critical sphere:
\begin{equation}
 \xi^T\nabla_q^2\RePart\lambda_*(q_0)\xi
 =\partial_r^2\RePart\lambda_*(q_0)
 (\xi\cdot\widehat q_0)^2.
 \label{eq:isotropic-tangential-degeneracy}
\end{equation}
Consequently full stability in $d\geq2$ additionally requires the angular
amplitude problem or the higher-order transverse phase equation.  The
present coefficients support, but do not yet prove, a stable isotropic
plane or hyperspace pattern.  An anisotropic isolated maximum can instead be
decided directly by \eqref{eq:sideband-J}--
\eqref{eq:sideband-diffusive-decay} when its diffusion tensor has strictly
damping real part.

\subsection{Extension 3: one periodic and several non-periodic dimensions}
\label{sec:mixed-domain}

Let
\begin{equation}
 (\theta,x)\in\mathbb S^1\times\mathbb R^d,\qquad
 0\leq\theta<2\pi,\qquad x\in\mathbb R^d,
\end{equation}
and let $\ell_\theta$ convert angular separation to spatial units.  With an
orientation width $\sigma_{ij}^\theta$ and spatial covariance $C_{ij}^x$,
the normalized product kernel has symbol
\begin{align}
 e_{n,k}(\theta,x)&=\ee^{\ii(n\theta+k\cdot x)},
 &n&\in\mathbb Z,\quad k\in\mathbb R^d,\notag\\
 \widehat J(n,k)_{ij}
 &=\mu_{ij}\exp\!\left[-\frac12\left(
 \frac{(\sigma_{ij}^\theta)^2n^2}{\ell_\theta^2}
 +k^TC_{ij}^xk\right)\right].
 \label{eq:mixed-symbol}
\end{align}
Using one width matrix $\Sigma$ means
$\sigma_{ij}^\theta=\sigma_{ij}$ and
$C_{ij}^x=\sigma_{ij}^2I_d$.  The modal matrix and exact spectral condition
are
\begin{align}
 M(n,k)&=\T^{-1}\left[-I+D_1\widehat J(n,k)\right],
 \label{eq:mixed-M}\\
 \sup_{n\in\mathbb Z,\,k\in\mathbb R^d}
 \max_s\RePart\lambda_s(M(n,k))&<0.
 \label{eq:mixed-stability}
\end{align}
The corresponding operator decomposition is
\begin{equation}
 \mathcal F_x\mathcal L\mathcal F_x^{-1}
 =\bigoplus_{n\in\mathbb Z}
 \int_{\mathbb R^d}^{\oplus}M(n,k)\,dk,
 \qquad
 \operatorname{spec}(\mathcal L)
 =\overline{\bigcup_{n\in\mathbb Z,\,k\in\mathbb R^d}
 \operatorname{spec}(M(n,k))}.
 \label{eq:mixed-direct-integral}
\end{equation}
Assume a finite critical set $(n_\alpha,k_\alpha)$ with simple smooth
eigenvalue bands.  After removing the group velocity in the non-periodic
directions, the leading expansion is
\begin{equation}
 u_i(\theta,x,t)=u_i^*+2\varepsilon\RePart\sum_{\alpha=1}^m
 A_\alpha(X,T)\Phi_{i\alpha}
 \ee^{\ii(n_\alpha\theta+k_\alpha\cdot x+\omega_\alpha t)}
 +O(\varepsilon^2),
 \qquad X=\varepsilon x,\quad T=\varepsilon^2t.
 \label{eq:mixed-superposition}
\end{equation}
The quadratic selection equations are
\begin{equation}
 n_r+n_s=n_\alpha,\qquad
 k_r+k_s=k_\alpha,\qquad
 \omega_r+\omega_s=\omega_\alpha,
 \label{eq:mixed-quadratic-resonance}
\end{equation}
and the cubic selection equations are
\begin{equation}
 n_\beta+n_\gamma-n_\delta=n_\alpha,\qquad
 k_\beta+k_\gamma-k_\delta=k_\alpha,\qquad
 \omega_\beta+\omega_\gamma-\omega_\delta=\omega_\alpha.
 \label{eq:mixed-resonance}
\end{equation}
Substitution of $(n,k)$ for $q$ in
\eqref{eq:general-slaved-mode}--\eqref{eq:general-resonant-amplitudes} gives
all local mixed-domain coefficients.  Denote that local polynomial by
$\mathcal N_\alpha^{\rm mix}(A,\overline A)$ and let
\begin{equation}
 H_\alpha^x=
 \nabla_k^2\lambda_\alpha(n_\alpha,k_\alpha).
 \label{eq:mixed-Hessian}
\end{equation}
The non-periodic directions add the envelope derivatives
\begin{equation}
 \partial_TA_\alpha
 =\mathcal N_\alpha^{\rm mix}(A,\overline A)
 -\frac12\sum_{r,s=1}^d(H_\alpha^x)_{rs}
 \partial_{X_rX_s}A_\alpha.
 \label{eq:mixed-envelope}
\end{equation}
Constant-envelope equilibria are the roots of
$\mathcal N^{\rm mix}=0$.  Their full stability is obtained from
\eqref{eq:sideband-J}--\eqref{eq:sideband-stability} with $H_\alpha$
replaced by $H_\alpha^x$.  This separates the discrete orientation
selection rule from continuous cortical sidebands.

The mixed domain has the same onset conditions as
\eqref{eq:finite-wave-threshold}--\eqref{eq:finite-wave-nondegeneracy} after
replacing $q$ by $(n,k)$ and differentiating only in the continuous variable
$k$.  The critical label must satisfy $(n_0,k_0)\ne(0,0)$, while the
homogeneous sector $(0,0)$ remains stable.  If orientation is axial,
$\theta\equiv\theta+\pi$, only even $n$ occur on the $2\pi$ circle.

\section{A three-population V1 example with broad excitation}

\subsection{What is measured and what is calibrated}

Billeh et al. constructed a data-driven mouse V1 model with Gaussian
distance-dependent connection probabilities.  Their kernel convention was
$\exp(-r^2/s^2)$, so the standard deviation in \eqref{eq:gaussian} is
$s/\sqrt2$.  The reported or assumed local widths used here are
\begin{center}
\begin{tabular}{@{}lrr@{}}
\toprule
Connection & $s$ ($\mu$m) & standard deviation $s/\sqrt2$ ($\mu$m)\\
\midrule
E $\to$ E   & 114 & 80.6\\
E $\to$ PV  & 92  & 65.1\\
E $\to$ SST & 103 & 72.8\\
PV $\to$ E  & 95  & 67.2\\
SST $\to$ E & 85  & 60.1\\
inhibitory $\to$ inhibitory & 120 & 84.9\\
\bottomrule
\end{tabular}
\end{center}
Thus E--E is broader than direct PV--E and SST--E in this local mouse-V1
dataset.  This supports the concern that the broad-inhibition assumption of a
two-population lateral-inhibition model should not be read literally.

The population weights below are dimensionless Wilson--Cowan gains calibrated
to put the measured width ratios near a finite-wave Hopf point.  They are
\emph{not} converted synaptic conductances:
\begin{equation}
 W=\begin{pmatrix}
 22.1035&-14.3515&-10.3472\\
 18.1160& -8.9407&-12.8206\\
  9.4942& -2.5978& -7.3034
 \end{pmatrix},
 \label{eq:Wnumeric}
\end{equation}
with rows $(E,PV,SST)$ as targets and columns $(E,PV,SST)$ as sources.  The
time constants are $(10,12,35)$ ms, the sigmoid equilibrium is
$\overline u=(0.2,0.2,0.2)$, and the cortical patch is $4.760$ mm long.
Biases co-vary with the common coupling scale so this equilibrium stays fixed.

The first crossing occurs at coupling scale $1.0000165$ and Fourier mode 3:
\begin{equation}
 q_0=0.39598\,(0.1\,\mathrm{mm})^{-1},\qquad
 \omega_0=0.089403\,\mathrm{ms}^{-1},\qquad
 f_0=14.2289\,\mathrm{Hz}.
\end{equation}
The cortical wavelength is $4.760/3=1.587$ mm.  At the same point, the
homogeneous mode has growth rate $-0.001695$ ms$^{-1}$, so the instability is
genuinely spatial rather than a homogeneous Hopf instability.

\begin{figure}[ht]
\centering
\includegraphics[width=\textwidth]{output/simulations/linear_spectrum.png}
\caption{Leading eigenvalue across spatial modes.  Mode 3 crosses first while
mode 0 remains damped.}
\label{fig:spectrum}
\end{figure}

The coefficient calculation from
\eqref{eq:bformula}--\eqref{eq:alphaformula} gives
\eqref{eq:numerical-coefficients}.  In particular, the broad E--E connection
does not prevent pattern formation once PV and SST are dynamically distinct.
The mechanism uses the full three-dimensional eigenstructure and cannot be
reduced to the sign of one difference of Gaussians.

\subsection{What biological data justify}

Three levels of evidence establish that V1 is more than a simple E/I split.
\begin{enumerate}[label=\arabic*.]
 \item Tasic et al. identified 111 transcriptomic types in mouse VISp.  Gouwens
 et al. measured 1,938 visual-cortical neurons and identified 17
 electrophysiological, 38 morphological, and 46 combined morpho-electric
 types.  Billeh et al. retained 17 coarse classes in a model of about 230,000
 mouse-V1 neurons.
 \item Adult mouse visual-cortex recordings show cell-type-specific connection
 motifs among morphologically defined interneurons.  Large-scale mouse and
 human recordings further show that inhibitory synaptic dynamics align with
 the presynaptic subclass.  Collapsing PV, SST, and VIP/Htr3a into one
 inhibitory variable discards exactly the differences that can change a Hopf
 instability.
 \item Horizontal axon collaterals of layer 2/3 pyramidal neurons extend for
 millimeters in tree-shrew V1 and become feature-specific outside roughly
 $500\,\mu$m.  Therefore broad and anisotropic excitation is anatomically
 plausible, especially in species with orientation maps.  A single isotropic
 Gaussian is already a major simplification.
\end{enumerate}

The rigorous conclusion is limited but useful: a model with broad excitation,
several inhibitory time scales, and cell-class-specific connections can have
a stable homogeneous state and a finite-wave oscillatory instability with
plausible cortical spatial and temporal scales.  The present data do not
identify the Wilson--Cowan gain matrix uniquely, and they do not establish
that three populations are necessary.  Necessity would require ruling out all
two-population kernels and nonlinearities, which the current evidence cannot
do.

\section{Cell-type and connection audit, then a five-population test}
\label{sec:cell-audit}

\subsection{The word ``all'' depends on measurement resolution}

The finest transcriptomic inventory and the finest complete dynamical model
are not the same object.  The Allen VISp/ALM SMART-seq taxonomy contains 133
types: 56 glutamatergic, 61 GABAergic, and 16 non-neuronal types
\cite{tasic2018}.  These identities do not come with a measured all-pairs
physiological connection matrix.  At the next coarser scale, Patch-seq and
anatomical work distinguish excitatory IT, ET, NP, and CT families and
inhibitory Pvalb, Sst, Vip, Lamp5, and Sncg families.  Still finer names can be
useful for targeting, but they cannot be inserted into \eqref{eq:field} until
their density, time scale, gain, and source--target coupling are known.

The Billeh model is presently the useful complete intermediate scale.  Its 17
populations and cell counts are
\begin{center}
\begin{tabular}{@{}lrrrr@{}}
\toprule
Layer & excitatory & Pvalb & Sst & Htr3a\\
\midrule
L1   & --     & --   & --   & 999\\
L2/3 & 12,689 & 640  & 464  & 1,107\\
L4   & 10,254 & 963  & 553  & 270\\
L5   & 7,569  & 613  & 555  & 117\\
L6   & 12,882 & 1,052& 1,059& 192\\
\bottomrule
\end{tabular}
\end{center}
Its published $17\times17$ table reports connection probability at a planar
separation of $75\,\mu$m.  The machine-readable transcription is
\texttt{data/billeh\_connection\_probabilities\_75um.csv}.  A blank remains
``undetermined''; it is not converted to a measured zero.  Billeh et al.
reported that 89 of 289 entries were undetermined and set them to zero when
instantiating the circuit \cite{billeh2020}.  Probability alone is not the
field weight $W_{ij}$: unitary PSP, convergence, population density, firing
gain, delay, and synaptic dynamics must also be combined.

Several reliable motifs survive across resolutions.  Pvalb cells provide
fast perisomatic inhibition and strong Pvalb--Pvalb coupling.  Sst cells act
more slowly on distal dendrites, avoid some within-Sst connections, and
inhibit several other classes.  Vip/Htr3a cells include disinhibitory cells
that preferentially target Sst.  Excitatory connections depend on distance
and functional similarity.  The MICrONS volume adds more than 0.5 billion
synapses across more than 200,000 cells; a proofread column resolved 1,352
neurons and more than 70,000 inhibitory synapses, including 18 connectivity
motif groups and widespread interneuron target specificity
\cite{microns2025,schneidermizell2025}.  This evidence argues for testing
several resolutions, not for claiming a completed fine-type matrix.

\subsection{Application after splitting SST-GER from local SST}

Drinnenberg et al. found that a median $66.1\%$ of labeled deep-L2/3 SST cells
met their general ensemble-response criterion, while $33.9\%$ did not
\cite{drinnenberg2025}.  The following five-population test preserves that
split.  Rows are targets and columns are sources in the order
$(E,PV,S_G,S_L,VIP)$:
{\small
\begin{equation}
W_5=\begin{pmatrix}
22.1035&-14.3515&-6.8395&-3.5077&-1.2000\\
18.1160& -8.9407&-8.4744&-4.3462&-1.5000\\
12.5000& -2.5978&-4.8276&-2.4759&-6.0000\\
 7.5000& -2.5978&-4.8276&-2.4759&-6.0000\\
 8.0000& -2.0000&-3.3050&-1.6950&-1.0000
\end{pmatrix}.
\label{eq:Wfive}
\end{equation}
}
The Billeh width ratios are retained except for E input to SST-GER, whose
standard deviation is set to $175\,\mu$m, inside the $125$--$225\,\mu$m range
explored by the preprint's broad-integration model.  The time constants are
$(10,12,45,25,20)$ ms.  The SST split and width are data-constrained; the
five-population gains in \eqref{eq:Wfive} are calibrated hypotheses, not
measured conductances.

The first instability occurs at scale $1.03988$, mode 4, wavelength
$1.190$ mm, and frequency $14.038$ Hz.  The homogeneous growth rate is still
$-0.00719$ ms$^{-1}$.  Projection gives \eqref{eq:five-coefficients}, hence
$\RePart c_5<\RePart b_5<0$ and one traveling wave is locally selected.  For
$a=0.01$, $\eta_t=\eta_x=0.04$, and zero detuning, the full complex
equal-amplitude state has
$r=0.14949$, $\phi=172.622^\circ$ modulo $180^\circ$, and Jacobian
eigenvalues
\begin{equation}
(-0.17302,-0.08170,-0.08000,-0.07861),
\end{equation}
so it is asymptotically stable.

\begin{figure}[ht]
\centering
\includegraphics[width=0.49\textwidth]{output/simulations/five_population_unforced.png}
\includegraphics[width=0.49\textwidth]{output/simulations/five_population_combined.png}
\caption{Five-population field.  Left: above onset, one counter-propagating
mode dominates.  Right: below onset, combined $2{:}1$ forcing makes the two
critical amplitudes nearly equal and pins the response.}
\label{fig:five-field}
\end{figure}

\section{Hallucination patterns explained by the model}

Let retinal polar coordinates be $(r,\theta)$ and approximate the
retino-cortical map by
\begin{equation}
 X=\log(r/r_0),\qquad Y=\theta.
 \label{eq:logmap}
\end{equation}
A cortical plane wave becomes the visual-field image
\begin{equation}
 P(r,\theta)=\cos\!\left(k_X\log(r/r_0)+k_Y\theta+\varphi\right).
 \label{eq:visual-pattern}
\end{equation}
Consequently
\begin{center}
\begin{tabular}{@{}L{0.24\textwidth}L{0.26\textwidth}L{0.38\textwidth}@{}}
\toprule
Cortical planform & Visual-field form & Common report\\
\midrule
$\cos(kX)$ & concentric rings & target, tunnel, funnel\\
$\cos(nY)$ & radial spokes & pinwheel, wheel, rays, cobweb-like spokes\\
$\cos(kX+nY)$ & logarithmic spiral & spiral, rotating or rocking form\\
three waves at $120^\circ$ & distorted lattice & honeycomb, grid, checkerboard\\
irregular band-limited waves & mottled or labyrinthine form & ripples, lines,
zigzags, TV-like structure\\
\bottomrule
\end{tabular}
\end{center}

\begin{figure}[ht]
\centering
\includegraphics[width=\textwidth]{output/simulations/retinotopic_hallucination_gallery.png}
\caption{Static model percepts obtained by applying the inverse logarithmic
retino-cortical map.  These are explanatory images, not flickering stimuli.}
\label{fig:gallery}
\end{figure}

These classes are observed outside mathematics.  Allefeld et al. exposed 20
participants to a spatially uniform $1$--$50$ Hz Ganzfeld scan.  They retained
168 recurring percepts and classified lines, ripples, tunnels, honeycombs,
wheels, rasters, and spirals; individual percept frequencies were often stable
across repeated scans, but frequency did not determine a unique percept.
The multistability of symmetry-related amplitude states is consistent with
that ambiguity.  A 2023 experiment with 20 screened participants compared
3, 8, 10, and 18 Hz rhythmic and arrhythmic stimulation and found the
strongest geometric forms and dynamics at 10 Hz rhythmic flicker.  The 2011
model produced primarily stripes at high forcing frequency and hexagons at
low forcing frequency, with two frequency islands.  These comparisons support
the model class, but they do not prove that every reported pattern originates
in V1.  Color, complex imagery, retinal phosphenes, thalamic effects, and
higher visual areas require additional mechanisms.

\section{Numerical implementation and results}

The file \texttt{neural\_field.py} accepts any number of populations, a signed
connectivity matrix, a Gaussian width for every source--target pair, time
constants, sigmoid gains, a homogeneous rest state, dimension, boundary
condition, and weak forcing.  It provides:
\begin{itemize}
 \item 1D and 2D simulations;
 \item periodic, reflecting, and zero exterior boundaries;
 \item direct $1{:}1$, temporal $2{:}1$, spatial $2{:}1$, and combined weak
 forcing;
 \item automatic critical-mode and Gaussian normal-form coefficients;
 \item Floquet exponents of every spatial mode about a forced uniform orbit;
 \item exact four-real-dimensional amplitude integration and Jacobians;
 \item static retino-cortical inputs and numerical temporal-drive arrays.
\end{itemize}
Periodic convolutions use Fourier multipliers
$\exp(-\sigma^2|q|^2/2)$.  Reflecting and zero-boundary simulations use
normalized Gaussian filtering in physical space.  Euler time stepping is
used deliberately so the code remains transparent; $\Delta t=0.5$ ms is
small relative to the shortest 10 ms population time constant.

\begin{figure}[ht]
\centering
\includegraphics[width=0.98\textwidth]{output/simulations/1d_periodic_unforced.png}
\includegraphics[width=0.98\textwidth]{output/simulations/1d_periodic_combined.png}
\caption{Top: above onset, one traveling mode dominates, as predicted by
$\RePart c<\RePart b$.  Bottom: below onset, combined temporal and spatial
$2{:}1$ forcing produces a standing, spatially pinned response.}
\label{fig:fields}
\end{figure}

The periodic unforced run above onset has late projected amplitudes
$\langle|A|\rangle=0.0103$ and $\langle|B|\rangle=0.1838$, so one direction
dominates.  With spatial $2{:}1$ forcing below onset, the two means are
$0.0800$ and $0.0681$, much closer to the standing-wave diagonal.  Reflecting
boundaries reduce the late E standard deviation from $0.1229$ in the
periodic combined-forcing run to $0.0773$ and visibly pin the edge response.
The 2D runs form banded and labyrinthine finite-wavenumber activity rather
than a homogeneous oscillation.  Exact values for every run are in
\texttt{output/simulations/summary.json}.

\section{Model inputs and the model-predicted frequency range}
\label{sec:floquet-range}

The five-population intrinsic frequency is $f_0=14.038$ Hz.  Direct and
$2{:}1$ resonant guesses are therefore $14.038$ and $28.076$ Hz, but a guessed
center frequency is not a range.  To calculate the range from the full field,
let $\overline u_F(t)$ be the spatially uniform periodic solution under a
uniform sinusoidal input of frequency $F$.  A perturbation of spatial mode
$q$ satisfies
\begin{equation}
 \dot v_q=\T^{-1}\left[-I+D_1(t)\widehat W(q)\right]v_q,\qquad
 D_1(t)=\operatorname{diag}F_i'(\overline h_{F,i}(t)).
 \label{eq:floquet-variation}
\end{equation}
If $P_q(F)$ is the fundamental matrix after one forcing period $T_F$, define
\begin{equation}
 \gamma_q(F)=\frac1{T_F}\log\rho(P_q(F)).
 \label{eq:floquet-growth}
\end{equation}
The uniform orbit is unstable to a pattern exactly when
$\max_{q\ne0}\gamma_q(F)>0$.  This calculation includes the time-periodic
sigmoid gain, switches among competing spatial modes, and requires no fitted
normal-form forcing coefficient.

Figure~\ref{fig:floquet} uses the five-population field at $0.99$ of critical
coupling.  At model-input amplitude $0.025$, positive finite-mode growth
occurs quasi-statically below approximately $8.8$ Hz, over the resolved bands
$9.6$--$13.2$ Hz and a narrow neighborhood of $14.0$ Hz, and over two
$2{:}1$ tongues near $25.0$--$26.6$ Hz and $27.2$--$28.6$ Hz.  Boundaries are
grid estimates from $0.5$ Hz spacing and change with coupling, input
amplitude, sigmoid operating point, and the population parameters.

\begin{figure}[ht]
\centering
\includegraphics[width=\textwidth]{output/simulations/floquet_resonance_map.png}
\caption{Left: maximum finite-mode Floquet growth versus temporal input
frequency and model-input amplitude.  The black contour is zero growth.
Right: the amplitude-$0.025$ slice; red shading marks model pattern growth.}
\label{fig:floquet}
\end{figure}

The numerical input bundle contains static target and pinwheel arrays, their
PNG renderings, and numerical time series at $f_0$ and $2f_0$ in
\texttt{output/simulations/model\_inputs/}.  It contains no animation or
autoplay display.  The computed Floquet intervals are a \emph{model
pattern-forming range}; they are not a measured human hallucination threshold
and they are not a safety range.  The Epilepsy Foundation identifies roughly
$5$--$30$ Hz as particularly relevant for photosensitive seizures, overlapping
the computed tongues, while W3C guidance avoids more than three flashes per
second.  Therefore no high-contrast human-viewable flicker intended to provoke
hallucinations is generated here.  Static high-contrast stripes and radial
patterns can also be provocative for pattern-sensitive viewers, so the static
arrays are model inputs rather than validated exposure stimuli.  Any human
stimulation requires an approved
clinical or psychophysics protocol with screening, luminance and contrast
limits, field-size limits, emergency stopping, and informed consent.

\section{Limitations and next experimental tests}

\begin{enumerate}[label=\arabic*.]
 \item The Gaussian width ratios are data-informed; the population weights in
 \eqref{eq:Wnumeric} are calibrated.  Direct conversion from connection
 probability and unitary synaptic strength to a Wilson--Cowan gain requires
 cell density, convergence, release dynamics, operating point, and gain.
 \item Mouse V1 has a salt-and-pepper orientation organization, while the
 classical human hallucination theory uses orientation columns.  Human,
 primate, carnivore, and tree-shrew horizontal connectivity should be used
 when spatially precise human predictions are intended.
 \item The cubic Turing--Hopf normal form explains the local high-frequency
 stripe selection.  Low-frequency hexagons require the Floquet normal form
 about the periodically forced homogeneous orbit and its quadratic term.
 \item The logarithmic map predicts geometry, not color.  Color opponency needs
 additional cone-opponent or cortical color populations.
 \item The strongest test is quantitative: measure each participant's damped
 visual frequency, then predict both the two forcing islands and the spatial
 wavelength before collecting hallucination reports.
\end{enumerate}

\clearpage
\appendix
\section{One-page summary: large-scale cellular-resolution read/write}
\label{app:drinnenberg}
\small
\textbf{Paper and question.}
Drinnenberg et al., ``Large-scale cellular-resolution read/write of activity
enables discovery of cell types defined by complex circuit properties,'' is a
2025 bioRxiv preprint and has not been peer reviewed
\cite{drinnenberg2025}.  The paper asks whether a cell type can be defined by
how it participates in a large causal circuit experiment, rather than only by
genes, morphology, or intrinsic physiology.  The central logic is to stimulate
many precisely chosen excitatory ensembles, read thousands of non-targeted
neurons, identify followers common to otherwise non-overlapping ensembles,
then use cell labels and a probabilistic model to explain that response class.

\textbf{Read/write experiment.}
New transgenic and follower-focused preparations combined soma-enriched
ChRmine with calcium indicators.  The platform controlled approximately
$10^3$ cells and monitored up to $10^4$ cells within or across mouse visual
areas.  In the main V1 experiment, 104 random, orientation-selective, or
scene-selective L2/3 excitatory ensembles were tested in 10 awake mice.  Each
ensemble had a median 50 targets and was driven for $1.5$ s at $32$ Hz while a
median 8,318 non-targeted regions of interest were monitored.  Downstream
activity decoded ensemble identity above chance; iso-tuned cells were
recruited while orthogonally tuned cells were suppressed, consistent with
like-to-like wiring.

\textbf{Probabilistic model.}
For a candidate follower, the connection probability from a pyramidal cell at
distance $r$ was
\begin{equation}
 p(r)=A\exp[-r^2/(2\sigma^2)].
\end{equation}
If $p_c$ is this probability averaged over uniformly distributed targets, the
probability of following an ensemble of $T$ targets with activation threshold
$k$ is
\begin{equation}
 P_f=1-\sum_{j=0}^{k-1}\binom{T}{j}p_c^j(1-p_c)^{T-j},
 \qquad \mathbb E[m]=M P_f
\end{equation}
for $M$ independently sampled ensembles.  With $M=10$, $T=50$, peak
connectivity $A=0.2$, $\sigma=100\,\mu$m, and $k=5$, the expected number
followed was only $0.004$ of 10.  Larger $\sigma$ helped, but a low $k$ was
also necessary.  The model therefore predicted a broadly integrating,
low-threshold inhibitory subtype rather than an ordinary local follower.

\textbf{Results.}
Single ensembles recruited $199\pm9$ positive followers ($2.5\%$ of the
non-targeted population) and suppressed $267\pm14$ cells ($3.4\%$).  Of the
positive followers, $18\pm3.8\%$ followed at least four ensembles, over
20-fold the chance expectation $0.87\pm0.3\%$; $6.3\pm1.8\%$ followed at
least eight.  These general ensemble-response (GER) neurons remained active
more than $150\,\mu$m from the nearest target, were concentrated in deep
L2/3, were broadly visually responsive, and were mutually correlated.
Cell-type labeling showed PV neurons were suppressed, whereas SST neurons were
increased.  A median $66.1\%$ of labeled SST cells met the GER criterion and
$33.9\%$ did not; $53.5\%$ of all functionally identified GER cells carried
the incomplete SST label, so the true SST fraction may be higher.

\textbf{Connection to this model.}
The result supports splitting SST-GER from local SST: the former receives a
broader E kernel, has a lower effective threshold and longer integration time,
and supplies dendritic feedback inhibition.  It does not provide a measured
five-by-five gain matrix.  Equation~\eqref{eq:Wfive} is therefore a transparent
hypothesis test.  Its finite-wave Hopf point and coefficients show that the
functional split is sufficient to preserve pattern formation with broad
excitation; they do not establish SST-GER dysfunction as a human
hallucination mechanism.
\normalsize
\clearpage

\begin{thebibliography}{99}

\bibitem{rule2011}
M. Rule, M. Stoffregen, and B. Ermentrout,
``A Model for the Origin and Properties of Flicker-Induced Geometric
Phosphenes,'' \emph{PLoS Computational Biology} 7, e1002158 (2011).
\href{https://doi.org/10.1371/journal.pcbi.1002158}{doi:10.1371/journal.pcbi.1002158}.

\bibitem{bressloff2002}
P. Bressloff, J. Cowan, M. Golubitsky, P. Thomas, and M. Wiener,
``What Geometric Visual Hallucinations Tell Us about the Visual Cortex,''
\emph{Neural Computation} 14, 473--491 (2002).
\href{https://doi.org/10.1162/089976602317250861}{doi:10.1162/089976602317250861}.

\bibitem{allefeld2011}
C. Allefeld, P. Putz, K. Kastner, and J. Wackermann,
``Flicker-light induced visual phenomena: Frequency dependence and
specificity of whole percepts and percept features,''
\emph{Consciousness and Cognition} 20, 1344--1362 (2011).
\href{https://doi.org/10.1016/j.concog.2010.10.026}{doi:10.1016/j.concog.2010.10.026}.

\bibitem{bartossek2023}
M. Bartossek et al.,
``Effect of frequency and rhythmicity on flicker light-induced hallucinatory
phenomena,'' \emph{PLoS ONE} 18, e0284271 (2023).
\href{https://doi.org/10.1371/journal.pone.0284271}{doi:10.1371/journal.pone.0284271}.

\bibitem{billock2007}
V. Billock and B. Tsou,
``Neural interactions between flicker-induced self-organized visual
hallucinations and physical stimuli,'' \emph{PNAS} 104, 8490--8495 (2007).
\href{https://doi.org/10.1073/pnas.0610813104}{doi:10.1073/pnas.0610813104}.

\bibitem{billeh2020}
Y. Billeh et al.,
``Systematic Integration of Structural and Functional Data into Multi-scale
Models of Mouse Primary Visual Cortex,'' \emph{Neuron} 106, 388--403 (2020).
\href{https://doi.org/10.1016/j.neuron.2020.01.040}{doi:10.1016/j.neuron.2020.01.040}.

\bibitem{tasic2018}
B. Tasic et al.,
``Shared and distinct transcriptomic cell types across neocortical areas,''
\emph{Nature} 563, 72--78 (2018).
\href{https://doi.org/10.1038/s41586-018-0654-5}{doi:10.1038/s41586-018-0654-5}.

\bibitem{gouwens2019}
N. Gouwens et al.,
``Classification of electrophysiological and morphological neuron types in
the mouse visual cortex,'' \emph{Nature Neuroscience} 22, 1182--1195 (2019).
\href{https://doi.org/10.1038/s41593-019-0417-0}{doi:10.1038/s41593-019-0417-0}.

\bibitem{jiang2015}
X. Jiang et al.,
``Principles of connectivity among morphologically defined cell types in
adult neocortex,'' \emph{Science} 350, aac9462 (2015).
\href{https://doi.org/10.1126/science.aac9462}{doi:10.1126/science.aac9462}.

\bibitem{campagnola2022}
L. Campagnola et al.,
``Local Connectivity and Synaptic Dynamics in Mouse and Human Neocortex,''
\emph{Science} 375, eabj5861 (2022).
\href{https://doi.org/10.1126/science.abj5861}{doi:10.1126/science.abj5861}.

\bibitem{bosking1997}
W. Bosking, Y. Zhang, B. Schofield, and D. Fitzpatrick,
``Orientation Selectivity and the Arrangement of Horizontal Connections in
Tree Shrew Striate Cortex,'' \emph{Journal of Neuroscience} 17, 2112--2127
(1997).
\href{https://doi.org/10.1523/JNEUROSCI.17-06-02112.1997}{doi:10.1523/JNEUROSCI.17-06-02112.1997}.

\bibitem{microns2025}
MICrONS Consortium,
``Functional connectomics spanning multiple areas of mouse visual cortex,''
\emph{Nature} 640, 435--447 (2025).
\href{https://doi.org/10.1038/s41586-025-08790-w}{doi:10.1038/s41586-025-08790-w}.

\bibitem{schneidermizell2025}
C. Schneider-Mizell et al.,
``Inhibitory specificity from a connectomic census of mouse visual cortex,''
\emph{Nature} 640, 448--458 (2025).
\href{https://doi.org/10.1038/s41586-024-07780-8}{doi:10.1038/s41586-024-07780-8}.

\bibitem{drinnenberg2025}
A. Drinnenberg et al.,
``Large-scale cellular-resolution read/write of activity enables discovery
of cell types defined by complex circuit properties,'' bioRxiv (2025).
\href{https://doi.org/10.1101/2025.10.21.683734}{doi:10.1101/2025.10.21.683734}.

\bibitem{photosensitivity}
Epilepsy Foundation,
``Photosensitivity and Seizures.''
\href{https://go.epilepsy.com/what-is-epilepsy/seizure-triggers/photosensitivity}{epilepsy.com}.

\bibitem{wcag}
W3C Web Accessibility Initiative,
``Understanding Success Criterion 2.3.2: Three Flashes.''
\href{https://www.w3.org/WAI/WCAG21/Understanding/three-flashes}{w3.org}.

\end{thebibliography}

\end{document}
